%=======================================================================
%  Chapter 2.  BODIES, CONFIGURATIONS, AND MOTIONS  (sections 2.1-2.5)
%  A from-scratch, motivated account of Steigmann & Shirani,
%  "Principles of Continuum Mechanics", Chapter 2, matching the book's
%  section and equation numbering.  Every step is shown, using only the
%  algebra of Chapter 1.  Examples are solved in place.
%
%  Compile:  pdflatex continuum_ch2.tex   (run twice, for references)
%=======================================================================
\documentclass[17pt,a4paper]{extreport}

\usepackage[margin=1.8cm]{geometry}
\usepackage{amsmath,amssymb,amsthm}
\usepackage{bm}
\usepackage{accents}
% book style: boldface is rendered by an under-tilde
\DeclareRobustCommand{\vek}[1]{\underaccent{\tilde}{#1}}
\usepackage{xcolor}
\usepackage{enumitem}
\usepackage{tikz}
\usetikzlibrary{arrows.meta,calc,positioning}

% reusable blob outline for bodies and configurations
\newcommand{\blob}[1]{\draw[#1] plot[smooth cycle,tension=0.85] coordinates
  {(0,1.05) (1.15,0.72) (1.5,-0.35) (0.6,-1.15) (-0.75,-0.95) (-1.3,0.05)};}
\tikzset{bodyline/.style={thick,dashed}, configline/.style={thick},
         vec/.style={-{Stealth[length=2.2mm]},thick},
         mapar/.style={-{Stealth[length=2.4mm]},thick}}
\usepackage[colorlinks=true,linkcolor=blue!55!black,citecolor=blue!55!black,
            hypertexnames=false,pdfencoding=unicode]{hyperref}

%=============================== notation ==============================
\newcommand{\Rn}{\mathbb{R}}
\newcommand{\Et}{\mathcal{E}^{3}}          % spatial point space
\newcommand{\hEt}{\hat{\mathcal{E}}^{3}}    % referential point space
\newcommand{\Vt}{\mathbb{E}^{3}}            % spatial translation space
\newcommand{\hVt}{\hat{\mathbb{E}}^{3}}     % referential translation space
\newcommand{\Lin}{\operatorname{Lin}}
\newcommand{\Proj}[1]{\mathbb{P}_{(#1)}}
\newcommand{\Refl}[1]{\mathbb{R}_{(#1)}}
\newcommand{\tr}{\operatorname{tr}}
\newcommand{\grad}{\operatorname{grad}}
\newcommand{\Grad}{\operatorname{Grad}}
\newcommand{\dvg}{\operatorname{div}}
\newcommand{\Dvg}{\operatorname{Div}}
\newcommand{\I}{\vek{I}}
\newcommand{\Ob}{\vek{O}}
\newcommand{\bE}{\vek{E}}
\newcommand{\ba}{\vek{a}}
\newcommand{\bb}{\vek{b}}
\newcommand{\bc}{\vek{c}}
\newcommand{\bd}{\vek{d}}
\newcommand{\bu}{\vek{u}}
\newcommand{\bv}{\vek{v}}
\newcommand{\bw}{\vek{w}}
\newcommand{\bx}{\vek{x}}
\newcommand{\by}{\vek{y}}
\newcommand{\bz}{\vek{z}}
\newcommand{\bX}{\vek{X}}
\newcommand{\bY}{\vek{Y}}
\newcommand{\bZ}{\vek{Z}}
\newcommand{\bo}{\vek{o}}
\newcommand{\be}{\vek{e}}
\newcommand{\bn}{\vek{n}}
\newcommand{\bN}{\vek{N}}
\newcommand{\bM}{\vek{M}}
\newcommand{\bmm}{\vek{m}}
\newcommand{\bze}{\vek{0}}
\newcommand{\bA}{\vek{A}}
\newcommand{\bB}{\vek{B}}
\newcommand{\bC}{\vek{C}}
\newcommand{\bF}{\vek{F}}
\newcommand{\bH}{\vek{H}}
\newcommand{\bL}{\vek{L}}
\newcommand{\bQ}{\vek{Q}}
\newcommand{\bR}{\vek{R}}
\newcommand{\bU}{\vek{U}}
\newcommand{\bV}{\vek{V}}
\newcommand{\hbE}{\hat{\vek{E}}}
\newcommand{\bchi}{\vek{\chi}}
\newcommand{\bkappa}{\vek{\kappa}}
\newcommand{\dd}{\delta}
\newcommand{\eps}{\varepsilon}
\newcommand{\dl}{\mathrm{d}}
\newcommand{\tbox}[3]{[\,#1,#2,#3\,]}
\newcommand{\ip}[2]{#1\cdot#2}
\newcommand{\rest}[1]{\big|_{#1}}

\newtheoremstyle{scplain}{\topsep}{\topsep}{\itshape}{0pt}
  {\scshape}{.}{5pt plus 1pt minus 1pt}{}
\newtheoremstyle{scdefn}{\topsep}{\topsep}{\normalfont}{0pt}
  {\scshape}{.}{5pt plus 1pt minus 1pt}{}

\theoremstyle{scplain}
\newtheorem{theorem}{Theorem}[chapter]
\newtheorem{proposition}[theorem]{Proposition}
\newtheorem{lemma}[theorem]{Lemma}
\newtheorem{corollary}[theorem]{Corollary}
\theoremstyle{scdefn}
\newtheorem{definition}[theorem]{Definition}
\newtheorem{remark}[theorem]{Remark}
\newtheorem{example}[theorem]{Example}

\renewcommand{\qedsymbol}{$\blacksquare$}
\allowdisplaybreaks
\setcounter{chapter}{1}   % so \chapter becomes Chapter 2

%=======================================================================

%------------------- clickable equation references ---------------------
% \eqr{4.11} prints "(4.11)" exactly as before and turns it into a link
% when equation (4.11) is present in this document.  When it is not -- a
% standalone chapter referring to an equation of another chapter -- it
% degrades silently to plain text instead of raising an undefined
% reference.  Must sit after hyperref.
\makeatletter
% a display carrying \tag{2.48--2.50} can hold only one \label, so the
% remaining members of the range are routed to it by this alias table
\newcommand{\eqalias}[2]{\expandafter\gdef\csname eqal@#1\endcsname{#2}}
\newcommand{\eqr@do}[1]{%
  \begingroup
    \@ifundefined{eqal@#1}{\def\eqr@t{#1}}%
                          {\edef\eqr@t{\csname eqal@#1\endcsname}}%
    \@ifundefined{r@eq:\eqr@t}%
      {\textup{(#1)}}%
      {\hyperref[eq:\eqr@t]{\textup{(#1)}}}%
  \endgroup}
\DeclareRobustCommand{\eqr}[1]{%
  \ifmmode\text{\eqr@do{#1}}\else\eqr@do{#1}\fi}
\makeatother
% generated: extra members of a range tag point at the tag itself
\eqalias{1.104}{1.103}
\eqalias{1.105}{1.103}
\eqalias{1.106}{1.103}
\eqalias{1.114}{1.113}
\eqalias{1.127}{1.126}
\eqalias{2.6}{2.5}
\eqalias{2.9}{2.8}
\eqalias{2.19}{2.18}
\eqalias{2.23}{2.22}
\eqalias{2.26}{2.25}
\eqalias{2.44}{2.43}
\eqalias{2.46}{2.45}
\eqalias{2.49}{2.48}
\eqalias{2.50}{2.48}
\eqalias{2.53}{2.52}
\eqalias{2.54}{2.52}
\eqalias{2.72}{2.71}
\eqalias{2.73}{2.71}
\eqalias{2.78}{2.77}
\eqalias{2.81}{2.80}
\eqalias{2.87}{2.86}
\eqalias{2.88}{2.86}
\eqalias{2.90}{2.89}
\eqalias{2.94}{2.93}
\eqalias{2.97}{2.96}
\eqalias{2.99}{2.98}
\eqalias{2.103}{2.102}
\eqalias{2.108}{2.107}
\eqalias{3.10}{3.9}


\begin{document}

\begin{titlepage}
\centering
\vspace*{3cm}
{\Large\bfseries Bodies, Configurations, and Motions\par}
\vspace{4mm}
{\large Kinematics of a deforming body: motion, the material derivative,\\
the deformation gradient, and mass\par}
\vspace{8mm}
{\itshape A from-scratch, definition--theorem--proof companion to}\\[2mm]
{Steigmann and Shirani, \emph{Principles of Continuum Mechanics},
Chapter~2}\\[1mm]
{(\S\,2.1 through \S\,2.5, in the book's order and numbering)}
\vfill
{\small Every figure of the text is described in an intuitive remark, every
equation is derived using only the algebra of Chapter~1, and the worked
example is solved in full. Section and equation numbers match the book.\par}
\vspace{1cm}
\end{titlepage}

\tableofcontents

\chapter{Bodies, Configurations, and Motions}

This chapter builds the \emph{kinematics} of continuum mechanics: the
language for describing where the material of a body is, how fast it moves,
and how it deforms, all before any mention of force. The single new idea is
that we carry \emph{two} pictures of the body at once, a fixed reference
copy and the moving current one, and translate between them. Everything
computational rests on the vector and tensor algebra of Chapter~1; we cite
it freely and prove each new relation from it.

%=======================================================================
\section{Position, velocity and acceleration}\label{sec:pva}
%=======================================================================

A \emph{body} $B$ is a collection of \emph{material points}, labelled $p$.
In the older literature these are sometimes called \emph{particles}, to
convey the notion of a small indivisible object; that name has some
heuristic value but is not really what is intended. A material point is
better regarded as a \emph{primitive} notion: not a tiny lump of matter, but
an abstract label for ``a bit of the material''. The body $B$ itself is not a
region of space; we never see $B$ directly, only the region of space it
fills at a given instant. In the book's figures $B$ is drawn with a
\emph{dashed} outline for exactly this reason (Figure~2.1): it is a set of
labels, not a place.

\begin{figure}[htbp]\centering
\begin{tikzpicture}[scale=0.85]
  \blob{bodyline}
  \fill (0.05,-0.05) circle (1.7pt) node[left=2pt] {$p$};
  \node at (1.45,0.95) {$B$};
\end{tikzpicture}
\caption{A material body.}\label{fig:2.1}
\end{figure}

\begin{figure}[htbp]\centering
\begin{tikzpicture}[scale=0.8]
  \begin{scope}\blob{bodyline}
    \fill (0,0) circle (1.7pt) node[left=2pt] {$p$};
    \node at (-1.15,1.0) {$B$};\end{scope}
  \begin{scope}[shift={(6,0)}]\blob{configline}
    \fill (0.1,0.1) circle (1.7pt) node[right=2pt] {$p$};
    \node at (1.3,1.0) {$\kappa_t$};\end{scope}
  \draw[mapar] (0.15,0.2) to[bend left=22] (5.9,0.35);
  \node at (3.1,1.5) {$\bx=\bchi(p,t)$};
\end{tikzpicture}
\caption{Configuration of the body at time $t$.}\label{fig:2.2}
\end{figure}

\begin{figure}[htbp]\centering
\begin{tikzpicture}[scale=0.8]
  \begin{scope}\blob{bodyline}
    \fill (0,0) circle (1.7pt) node[left=2pt] {$p$};
    \node at (-1.15,1.0) {$B$};\end{scope}
  \begin{scope}[shift={(6,0)}]\blob{configline}
    \fill (0.1,0.1) circle (1.7pt) node[right=2pt] {$p$};
    \node at (1.3,1.0) {$\kappa$};\end{scope}
  \draw[mapar] (0.15,0.2) to[bend left=22] (5.9,0.35);
  \node at (3.1,1.5) {$\bX=\bkappa(p)$};
\end{tikzpicture}
\caption{A reference configuration for the body.}\label{fig:2.3}
\end{figure}

\begin{remark}[Figures 2.1--2.3: body versus configuration]\label{rem:fig123}
Figure~2.1 shows the abstract body $B$ (a dashed blob with a marked point
$p$): the reservoir of material points. Figures~2.2 and~2.3 each show $B$ on
the left mapped by an arrow to a \emph{solid} blob on the right, the
difference being only which placement of the body the arrow produces: the
\emph{current} configuration $\kappa_t$ at time $t$ in Figure~2.2, and a
fixed \emph{reference} configuration $\kappa$ in Figure~2.3. Solid outlines
mark genuine regions of Euclidean space; the dashed $B$ is not one. The
whole chapter is the art of never confusing the label $p$ with either of
its placements $\bx$ or $\bX$.
\end{remark}

\begin{definition}[Configuration and motion]\label{def:motion}
At time $t$ the body occupies a \emph{configuration} $\kappa_t$, a region of
the Euclidean point space $\Et$, and $p\in B$ sits at the position
$\bx\in\kappa_t$. Distinct material points occupy distinct positions at a
fixed time, so there is an invertible \emph{motion} $\bchi(p,t)$ with
\begin{equation}
   \bx=\bchi(p,t),\qquad p=\bchi^{-1}(\bx,t). \tag{2.1}\label{eq:2.1}
\end{equation}
\end{definition}

At fixed $t$, $\bchi$ tells us where each material point is, and
$\bchi^{-1}$ tells us which material point sits at a given place. Freezing
$p$ and letting $t$ run traces the path of one material point through
space; its time derivatives are its velocity and acceleration.

\begin{definition}[Velocity, acceleration]\label{def:va}
\begin{equation}
   \bv(p,t)=\frac{\partial}{\partial t}\bchi(p,t),\qquad
   \ba(p,t)=\frac{\partial^{2}}{\partial t^{2}}\bchi(p,t). \tag{2.2}\label{eq:2.2}
\end{equation}
The derivatives are taken at fixed $p$: they follow the \emph{same} material
point.
\end{definition}

Working with the ethereal $B$ is awkward, because a ``label'' has no
coordinates. The remedy is to name each material point by the fixed place
it occupies in one chosen, frozen configuration, the \emph{reference}
configuration $\kappa$.

\begin{definition}[Reference configuration]\label{def:ref}
Fix a configuration $\kappa$ and let $\bX\in\kappa$ be positions in one-to-one
correspondence with material points through an invertible placement map
$\bkappa$,
\begin{equation}
   \bX=\bkappa(p),\qquad p=\bkappa^{-1}(\bX). \tag{2.3}\label{eq:2.3}
\end{equation}
A common, though not obligatory, choice is the configuration actually
occupied at some instant $t_0$, $\kappa=\kappa_{t_0}$, for which
\begin{equation}
   \bkappa(p)=\bchi(p,t_0). \tag{2.4}\label{eq:2.4}
\end{equation}
\end{definition}

Composing the placement's inverse with the motion expresses the current
position directly in terms of the reference position (Figure~2.4):
\begin{equation}
   \bx=\bchi_\kappa(\bX,t),\qquad
   \bchi_\kappa(\bX,t)=\bchi\!\big(\bkappa^{-1}(\bX),t\big). \tag{2.5--2.6}\label{eq:2.5}
\end{equation}
The subscript records the dependence on the chosen $\kappa$. Since each
$p$ corresponds to exactly one $\bx\in\kappa_t$ and one $\bX\in\kappa$, the
map $\bchi_\kappa(\cdot,t)$ is invertible at each $t$,
\begin{equation}
   \bX=\bchi_\kappa^{-1}(\bx,t). \tag{2.7}\label{eq:2.7}
\end{equation}

\begin{figure}[htbp]\centering
\begin{tikzpicture}[scale=0.82]
  \begin{scope}\blob{configline}
    \fill (0,0) circle (1.7pt) node[above right=0pt and 2pt] {$p$};
    \node at (-1.25,1.25) {$\kappa$};\end{scope}
  \fill (-2.1,-2.5) circle (1.4pt) node[below=2pt] {$O$};
  \draw[vec] (-2.1,-2.5) -- (-0.05,-0.06);
  \node at (-1.75,-1.2) {$\bX$};
  \begin{scope}[shift={(6.6,0)}]\blob{configline}
    \fill (0.15,0.1) circle (1.7pt) node[above right=0pt and 2pt] {$p$};
    \node at (1.45,1.25) {$\kappa_t$};\end{scope}
  \fill (4.5,-2.5) circle (1.4pt) node[below=2pt] {$O$};
  \draw[vec] (4.5,-2.5) -- (6.72,0.04);
  \node at (4.85,-1.2) {$\bx$};
  \draw[mapar] (0.75,0.95) to[bend left=20] (5.75,1.05);
  \node at (3.3,2.15) {$\bx=\bchi_\kappa(\bX,t)$};
\end{tikzpicture}
\caption{Mapping of reference configuration to current configuration.}\label{fig:2.4}
\end{figure}

\begin{remark}[Figure 2.4: the deformation map]\label{rem:fig4}
Figure~2.4 draws two \emph{solid} blobs, $\kappa$ (reference, left) and
$\kappa_t$ (current, right), with the same material point $p$ appearing in
each, at $\bX$ and at $\bx$. The connecting arrow is $\bchi_\kappa$: it is
the object a deformation actually is, a rule sending each reference place to
its current place. Because both blobs are real regions of space, this map is
the one we can differentiate, and its gradient will be the deformation
gradient of \S\ref{sec:defgrad}.
\end{remark}

Because $B$ and $\kappa$ are in fixed one-to-one correspondence, using
$(p,t)$ as independent variables is the same as using $(\bX,t)$. The
velocity therefore has a \emph{referential} form (a function of $\bX,t$) and
a \emph{spatial} form (a function of $\bx,t$),
\begin{equation}
   \hat\bv(\bX,t)=\bv\big(\bkappa^{-1}(\bX),t\big)
   =\frac{\partial}{\partial t}\bchi_\kappa(\bX,t),\qquad
   \tilde\bv(\bx,t)=\bv\big(\bchi^{-1}(\bx,t),t\big), \tag{2.8--2.9}\label{eq:2.8}
\end{equation}
the partial in \eqr{2.8} being taken at fixed $\bX$, i.e.\ at fixed $p$. The
two forms are related through \eqr{2.5} and \eqr{2.7},
\begin{equation}
   \hat\bv(\bX,t)=\tilde\bv(\bchi_\kappa(\bX,t),t),\qquad
   \tilde\bv(\bx,t)=\hat\bv(\bchi_\kappa^{-1}(\bx,t),t). \tag{2.10}\label{eq:2.10}
\end{equation}

%=======================================================================
\section{Material, referential and spatial descriptions}\label{sec:desc}
%=======================================================================

Any field $f$ carried by the material, temperature, density, velocity, may
be written as a function of the material point, of the current position, or
of the reference position:
\begin{equation}
   f=\bar f(p,t)=\tilde f(\bx,t)=\hat f(\bX,t), \tag{2.11}\label{eq:2.11}
\end{equation}
where
\begin{equation}
   \tilde f(\bx,t)=\bar f\big(\bchi^{-1}(\bx,t),t\big),\qquad
   \hat f(\bX,t)=\bar f\big(\bkappa^{-1}(\bX),t\big), \tag{2.12}\label{eq:2.12}
\end{equation}
and, relating the last two directly,
\begin{equation}
   \hat f(\bX,t)=\tilde f(\bchi_\kappa(\bX,t),t),\qquad
   \tilde f(\bx,t)=\hat f(\bchi_\kappa^{-1}(\bx,t),t). \tag{2.13}\label{eq:2.13}
\end{equation}
These are the \emph{material}, \emph{spatial}, and \emph{referential}
descriptions of $f$.

\begin{remark}[same value, three different functions: why the decorations
matter]\label{rem:decorations}
This is the single most important notational point of the chapter, and it is
the exact analogue of the caret in Chapter~1, where $\hat\varphi(x_1,x_2,x_3)$
and $\varphi(\bx)$ shared their values but had different domains
($\Rn^3$ versus a point space) and so were different functions. Here the
\emph{same number} $f$, the temperature of a given bit of matter at a given
instant, is produced by \emph{three different machines}:
\[
\begin{aligned}
&\bar f\ \text{eats a material label }p,\\
&\tilde f\ \text{eats a current place }\bx,\\
&\hat f\ \text{eats a reference place }\bX .
\end{aligned}
\]
They agree in value only when $p$, $\bx$, $\bX$ describe the same material
point at the same time, i.e.\ when $\bx=\bchi_\kappa(\bX,t)$. Their
\emph{derivatives} are genuinely different, because ``hold the variable
fixed'' means hold $p$, hold $\bx$, or hold $\bX$, and these are three
different physical instructions. The bar, tilde and caret are the honest
bookkeeping of which machine is meant; the book (and we) will soon drop the
decorations ``to minimize visual distraction'', exactly as the caret was
dropped in Chapter~1, on the understanding that the reader remembers which
function a symbol denotes.
\end{remark}

In the \emph{spatial} description $\bx$ and $t$ are the independent
variables. Fixing $\bx$ fixes a \emph{place}, not a material point, so as
time passes different material points stream through that fixed $\bx$
(Figure~2.5). The set of fixed positions $\bx$ is a \emph{control volume},
a window in space through which material flows.

\begin{figure}[htbp]\centering
\begin{tikzpicture}[scale=0.95]
  % same body drawn at two instants; it translates past a fixed place
  \begin{scope}\blob{configline}\end{scope}
  \node at (-1.5,1.35) {$\kappa_{t_1}$};
  \begin{scope}[shift={(1.65,0)}]\blob{thick,gray,dash pattern=on 3pt off 2pt}\end{scope}
  \node[gray] at (3.35,1.3) {$\kappa_{t_2}$};
  % the fixed place in space, marked by a cross
  \coordinate (X) at (0.55,-0.15);
  \draw[thin] ($(X)+(-0.32,0)$)--($(X)+(0.32,0)$);
  \draw[thin] ($(X)+(0,-0.32)$)--($(X)+(0,0.32)$);
  \fill (X) circle (2pt);
  \node[above left=2pt and 1pt] at (X) {$p_1$};
  % p1 is convected away while x stays put
  \draw[-{Stealth[length=2mm]},gray,densely dashed] ($(X)+(0.18,0)$) -- ++(1.3,0);
  \fill[gray] ($(X)+(1.65,0)$) circle (1.8pt);
  \node[gray,above=3pt] at ($(X)+(1.65,0)$) {$p_1$ at $t_2$};
  \node at (0.9,-1.95) {$\bx$ stays put; at $t_2$ a \emph{different} point $p_2$ sits there};
\end{tikzpicture}
\caption{Different material points occupy position $\bx$ at different times.}\label{fig:2.5}
\end{figure}

\begin{remark}[Figure 2.5: the control-volume viewpoint]\label{rem:fig5}
Figure~2.5 shows one fixed position $\bx$ occupied by different material
points at different times: the Eulerian, or spatial, standpoint of an
observer nailed to space (a thermometer on a wall) rather than riding with
the matter (a thermometer floating down a river). The distinction between
``rate of change at a fixed place'' and ``rate of change following the
matter'' is precisely what the material derivative below measures.
\end{remark}

Take the differential of $f$ in each description. Using the coordinate free
chain rule of Chapter~1, $\dl f=\ip{\grad f}{\dl\bx}$ (there written
$\dl\varphi=\ip{\grad\varphi}{\dl\bx}$), and allowing time to vary,
\begin{align}
   \dl f&=\tilde f'\,\dl t+\ip{\grad\tilde f}{\dl\bx}, \tag{2.14}\label{eq:2.14}\\
   \dl f&=(\hat f)^{\hbox{\tiny\textbullet}}\,\dl t+\ip{\Grad\hat f}{\dl\bX}. \tag{2.15}\label{eq:2.15}
\end{align}
Here $\tilde f'=\partial\tilde f/\partial t$ is the \emph{spatial time
derivative} (at fixed $\bx$) and $\grad$ is the gradient in $\bx$; while
$(\hat f)^{\hbox{\tiny\textbullet}}=\partial\hat f/\partial t$ is the
\emph{referential} (equivalently \emph{material}) time derivative (at fixed
$\bX$, i.e.\ fixed $p$) and $\Grad$ is the gradient in $\bX$. Since \eqr{2.14}
and \eqr{2.15} are two expressions for the \emph{same} increment $\dl f$,
dropping decorations,
\begin{equation}
   \dot f\,\dl t+\ip{\Grad f}{\dl\bX}=f'\,\dl t+\ip{\grad f}{\dl\bx}. \tag{2.16}\label{eq:2.16}
\end{equation}
We call $\dot f$ the \emph{material derivative}: the rate of change of $f$
following a fixed material point.

\begin{theorem}[Material derivative in the spatial description]\label{thm:matder}
For a scalar field, with $\bv$ the (spatial) material velocity,
\begin{equation}
   \dot f=f'+\ip{\grad f}{\bv}. \tag{2.17}\label{eq:2.17}
\end{equation}
For a vector field $\bu$,
\begin{equation}
   \dot\bu=\bu'+(\grad\bu)\bv, \tag{2.20}\label{eq:2.20}
\end{equation}
and in particular the material acceleration is
\begin{equation}
   \ba=\dot\bv=\bv'+\bL\bv,\qquad \bL=\grad\bv. \tag{2.21}\label{eq:2.21}
\end{equation}
\end{theorem}

\begin{proof}
Follow one material point: hold $p$, hence $\bX$, fixed, so $\dl\bX=\bze$ in
\eqr{2.16}. Moving with the point, $\dl\bx=\dl\bchi_\kappa(\bX,t)=\bv\,\dl t$
(the position changes at the velocity). Substituting into \eqr{2.16},
\[
   \dot f\,\dl t=f'\,\dl t+\ip{\grad f}{\bv\,\dl t}
   =\big(f'+\ip{\grad f}{\bv}\big)\dl t,
\]
and cancelling $\dl t$ gives \eqr{2.17}. For a vector field the identical
argument applied componentwise, or directly from the vector chain rule
$\dl\bu=\bu'\,\dl t+(\grad\bu)\dl\bx$ of Chapter~1, gives
\begin{equation}
   \dot\bu\,\dl t+(\Grad\bu)\dl\bX=\dl\bu=\bu'\,\dl t+(\grad\bu)\dl\bx,
   \tag{2.18--2.19}\label{eq:2.18}
\end{equation}
and setting $\dl\bX=\bze$, $\dl\bx=\bv\,\dl t$ yields \eqr{2.20}. Taking
$\bu=\bv$ gives the acceleration \eqr{2.21}, with $\bL=\grad\bv$ the spatial
velocity gradient.
\end{proof}

\begin{remark}[the convective term, intuitively]\label{rem:convective}
The material derivative splits into two clearly physical pieces:
\[
   \dot f=\underbrace{f'}_{\substack{\text{local: change at}\\ \text{a fixed place}}}
   +\underbrace{\ip{\grad f}{\bv}}_{\substack{\text{convective: change from}\\
   \text{being carried to new places}}}.
\]
Float down a river whose temperature is steady in time ($f'=0$) but rises
downstream ($\grad f\neq\bze$): you still warm up, at the rate
$\ip{\grad f}{\bv}$, purely because the current carries you into warmer
water. The acceleration \eqr{2.21} shows the same effect is why steady flows
($\bv'=\bze$) can still accelerate fluid particles, through the nonlinear
term $\bL\bv=(\grad\bv)\bv$; this convective nonlinearity is the source of
most of the difficulty in fluid and solid dynamics.
\end{remark}

\begin{figure}[htbp]\centering
\begin{tikzpicture}[scale=1.0]
  % channel
  \fill[gray!8] (0,0) rectangle (8.2,1.7);
  \draw[gray!60] (0,0)--(8.2,0); \draw[gray!60] (0,1.7)--(8.2,1.7);
  % isotherms: field increasing to the right
  \foreach \x/\lab in {1.0/{}, 2.6/{}, 4.2/{}, 5.8/{}, 7.4/{}}
     {\draw[gray!55,densely dashed] (\x,0)--(\x,1.7);}
  \node[gray!70] at (1.0,-0.3) {\scriptsize$f$ small};
  \node[gray!70] at (7.4,-0.3) {\scriptsize$f$ large};
  \draw[-{Stealth[length=2mm]},gray!70] (0.9,-0.62)--(7.5,-0.62);
  \node[gray!70] at (4.2,-0.95) {\scriptsize$\grad f$ points downstream};
  % flow velocity
  \draw[vec] (1.2,1.25)--(2.5,1.25) node[above=0pt,midway] {\small$\bv$};
  % fixed probe
  \draw[thick] (2.6,1.7)--(2.6,0.95); \fill (2.6,0.95) circle (2.2pt);
  \node[align=center,font=\scriptsize] at (1.75,2.35)
     {probe bolted to the bank:\\ reads $f'$ only};
  % floating probe, two instants
  \fill (4.6,0.62) circle (2.2pt); \node[font=\scriptsize,below=2pt] at (4.6,0.62) {$t$};
  \draw[-{Stealth[length=2mm]},densely dashed] (4.75,0.62)--(6.35,0.62);
  \fill (6.5,0.62) circle (2.2pt); \node[font=\scriptsize,below=2pt] at (6.5,0.62) {$t+\dl t$};
  \node[align=center,font=\scriptsize] at (6.2,2.35)
     {probe drifting with the water:\\ reads $\dot f=f'+\ip{\grad f}{\bv}$};
\end{tikzpicture}
\\[1.8ex]
{\footnotesize\textsc{Supplementary Figure A.} The two time derivatives of a
field. A probe bolted to the bank samples one \emph{place} and records only
the explicit time change $f'$. A probe drifting with the water samples one
\emph{material point}; even in a perfectly steady field ($f'=0$) it registers
a change, because the current carries it into regions of different $f$. That
extra reading is the convective term $\ip{\grad f}{\bv}$, and the sum is the
material derivative $\dot f$ of \eqr{2.17}.}
\end{figure}

%=======================================================================
\section{Cartesian coordinates}\label{sec:cart}
%=======================================================================

To compute, we resolve everything on orthonormal bases. The current
position uses a fixed basis $\{\be_i\}$ of the spatial translation space
$\Vt$, the reference position a second fixed basis $\{\hbE_A\}$ of the
referential translation space $\hVt$ (Figure~2.6):
\begin{equation}
   \bx=x_i\be_i\ (i\in\{1,2,3\}),\qquad
   \bX=X_A\hbE_A\ (A\in\{1,2,3\}). \tag{2.22--2.23}\label{eq:2.22}
\end{equation}

\begin{figure}[htbp]\centering
\begin{tikzpicture}[scale=0.82]
  % reference
  \begin{scope}[shift={(0.35,1.7)}]\blob{configline}
    \fill (0.2,0.05) circle (1.7pt) node[above right=0pt and 2pt] {$p$};\end{scope}
  \node at (0.35,3.35) {$\kappa\subset\hEt$};
  \coordinate (O1) at (-1.9,-1.4);
  \draw[vec] (O1) -- ++(1.25,0) node[right=1pt] {$\hbE_1$};
  \draw[vec] (O1) -- ++(0,1.15) node[above=1pt] {$\hbE_2$};
  \draw[vec] (O1) -- ++(-0.62,-0.58) node[below=0pt] {$\hbE_3$};
  \draw[vec] (O1) -- (0.5,1.7);
  \node at (-1.25,0.5) {$\bX$};
  % current
  \begin{scope}[shift={(7.0,1.7)}]\blob{configline}
    \fill (0.2,0.05) circle (1.7pt) node[above right=0pt and 2pt] {$p$};\end{scope}
  \node at (7.0,3.35) {$\kappa_t\subset\Et$};
  \coordinate (O2) at (4.75,-1.4);
  \draw[vec] (O2) -- ++(1.25,0) node[right=1pt] {$\be_1$};
  \draw[vec] (O2) -- ++(0,1.15) node[above=1pt] {$\be_2$};
  \draw[vec] (O2) -- ++(-0.62,-0.58) node[below=0pt] {$\be_3$};
  \draw[vec] (O2) -- (7.15,1.7);
  \node at (5.4,0.5) {$\bx$};
\end{tikzpicture}
\caption{Position vectors in $\kappa$ and $\kappa_t$.}\label{fig:2.6}
\end{figure}

\begin{remark}[Figure 2.6, and why two bases]\label{rem:fig6}
Figure~2.6 draws the reference blob $\kappa\subset\hEt$ with its own frame
$\{\hbE_A\}$ and the current blob $\kappa_t\subset\Et$ with a separate frame
$\{\be_i\}$. Keeping the two frames \emph{conceptually distinct}, even
though both are copies of ordinary three-space, is the whole trick of the
chapter: a quantity carrying one index of each kind (a lower-case $i$ from
$\{\be_i\}$ and an upper-case $A$ from $\{\hbE_A\}$) is a \emph{two-point
tensor}, an object living with one foot in each configuration. The
deformation gradient is the first and most important example. This is why
the book uses capital letters $X_A$, $\hbE_A$ for reference quantities and
lower-case $x_i$, $\be_i$ for current ones: the alphabet itself tracks which
space you are in.
\end{remark}

Regarding $\bchi_\kappa$ as a vector in $\Vt$ and using that the $\hbE_A$ are
fixed,
\begin{equation}
   x_i\be_i=\bchi_\kappa(X_A\hbE_A,t)=\chi_i(X_A,t)\,\be_i,\qquad
   x_i=\chi_i(X_A,t), \tag{2.24}\label{eq:2.24}
\end{equation}
shorthand for $x_i=\chi_i(X_1,X_2,X_3,t)$. The velocity, in the referential
description, is obtained by differencing at fixed $\bX$ and passing to the
limit:
\begin{equation}
\begin{aligned}
   \bv&=\lim_{\alpha\to0}\tfrac1\alpha\big[\bchi_\kappa(\bX,t+\alpha)
      -\bchi_\kappa(\bX,t)\big]\\
   &=\be_i\lim_{\alpha\to0}\tfrac1\alpha\big[\chi_i(X_A,t+\alpha)
      -\chi_i(X_A,t)\big]=\dot\chi_i\,\be_i,
\end{aligned}\tag{2.25--2.26}\label{eq:2.25}
\end{equation}
so $\hat v_i=\dot\chi_i=\partial\chi_i/\partial t$; likewise
$\hat a_i=\ddot\chi_i=\partial^2\chi_i/\partial t^2$ $(2.27)$. In the spatial
description $\bv=\tilde v_i(x_k,t)\be_i$, $\ba=\tilde a_i(x_k,t)\be_i$
$(2.28)$. Resolving the material derivative \eqr{2.21}, $\ba=\bv'+(\grad\bv)\bv$,
on $\{\be_i\}$ gives the component identity
\begin{equation}
   \tilde a_i=\tilde v_i'+\tilde v_{i,j}\tilde v_j, \tag{2.29}\label{eq:2.29}
\end{equation}
and applying \eqr{2.20} to the position field $\bx$ itself, with $\bx'=\bze$
and $x_{i,j}=\dd_{ij}$, recovers the sanity check $\dot x_i=x_{i,j}\tilde v_j
=\dd_{ij}\tilde v_j=\tilde v_i$: the material derivative of position is
velocity, as it must be.

\begin{remark}[why $\bx'=\bze$, and why this is not circular]\label{rem:xprime}
This step confuses almost everyone on first reading, because $\bx$ is
manifestly moving: the material point does change its position with time.
The resolution is that the prime is \emph{not} a material derivative.

Take the field $\bu$ in \eqr{2.18} to be \emph{position itself}. Its three
descriptions are
\[
   \bar\bu(p,t)=\bchi(p,t),\qquad
   \hat\bu(\bX,t)=\bchi_\kappa(\bX,t),\qquad
   \tilde\bu(\bx,t)=\bx .
\]
The first two genuinely depend on $t$: hold the material point fixed and its
location changes. The third does not, and the reason is purely definitional.
In the spatial description the independent variable \emph{is} $\bx$. Asking
for the value of the position field at the place $\bx$ is asking ``what
position is the position $\bx$?'', and the answer is $\bx$, at every instant.
The function $\tilde\bu$ is the identity map of $\Et$, and the identity map
carries no $t$ in it at all. Hence
\[
   \bu'=\frac{\partial}{\partial t}\Big|_{\bx\ \mathrm{fixed}}\tilde\bu(\bx,t)
       =\frac{\partial\bx}{\partial t}\Big|_{\bx\ \mathrm{fixed}}=\bze,
   \qquad\text{componentwise } \tilde u_i'=\frac{\partial x_i}{\partial t}=0 .
\]
Nothing is being asserted about the motion here. What is being asserted is
that a photograph of space, labelled by its own coordinates, does not change
when the clock advances. The matter flowing through the point changes; the
point does not.

The second ingredient is $\grad\bu=\I$, for the same reason: from
$\tilde u_i=x_i$ we get $u_{i,j}=\partial x_i/\partial x_j=\dd_{ij}$, so
$\grad\bu=\dd_{ij}\be_i\otimes\be_j=\I$. Substituting both into \eqr{2.20},
\[
   \dot\bu=\bu'+(\grad\bu)\bv=\bze+\I\bv=\bv,
\]
that is $\dot x_i=\tilde v_i$. Read backwards this is a consistency check on
the whole apparatus rather than a new result: \eqr{2.17} and \eqr{2.20} were
constructed precisely so that $\dot f=f'+\grad f\cdot\bv$ would reproduce the
rate seen by a material point, and the material point's rate of change of
position is, by \eqr{2.2}$_1$, its velocity. Had the formula returned anything
other than $\bv$, the definition of the material derivative would have been
wrong.

The whole content of the identity is therefore the split
\[
   \underbrace{\dot x_i}_{\text{the point moves}}
   =\underbrace{x_i'}_{=\,0:\ \text{the \emph{place} does not move}}
   +\underbrace{x_{i,j}\tilde v_j}_{=\,\tilde v_i:\ \text{convection}}.
\]
All of the motion sits in the convective term, none of it in the spatial
time derivative. This is the cleanest possible illustration of the difference
between $'$ and $\dot{\phantom{f}}$: for the position field the two are as
far apart as they can be, $\bze$ against $\bv$.
\end{remark}

\begin{example}[converting a motion between descriptions]\label{ex:motion}
Consider the motion (linear in $\bX$, hence atypical, but transparent)
\begin{equation}
   \begin{pmatrix}x_1\\x_2\\x_3\end{pmatrix}
   =\begin{pmatrix}e^{t}&0&0\\0&1&t\\0&-t&1\end{pmatrix}
   \begin{pmatrix}X_1\\X_2\\X_3\end{pmatrix},\quad\text{i.e.}\quad
   \begin{aligned}x_1&=e^{t}X_1,\\ x_2&=X_2+tX_3,\\ x_3&=-tX_2+X_3.\end{aligned}
   \tag{2.30}\label{eq:2.30}
\end{equation}

\emph{Referential velocity and acceleration.} Differentiate at fixed $\bX$:
\[
\begin{aligned}
   &\hat v_1=\dot\chi_1=e^{t}X_1,\quad \hat v_2=X_3,\quad \hat v_3=-X_2;\\
   &\hat a_1=e^{t}X_1,\quad \hat a_2=0,\quad \hat a_3=0 .
\end{aligned}
\]

\emph{Inversion.} To get the spatial forms we invert \eqr{2.30}. Immediately
$X_1=e^{-t}x_1$. The lower $2\times2$ block $\left(\begin{smallmatrix}1&t\\
-t&1\end{smallmatrix}\right)$ has determinant $1+t^2$ and inverse
$\tfrac{1}{1+t^2}\left(\begin{smallmatrix}1&-t\\t&1\end{smallmatrix}\right)$,
so
\begin{equation}
   X_1=e^{-t}x_1,\qquad
   X_2=\frac{x_2-t x_3}{1+t^2},\qquad
   X_3=\frac{t x_2+x_3}{1+t^2}. \tag{2.31}\label{eq:2.31}
\end{equation}

\emph{Spatial velocity.} Substitute \eqr{2.31} into $\hat v_i$:
\[
\begin{aligned}
   &\tilde v_1=e^{t}X_1=e^{t}e^{-t}x_1=x_1,\\
   &\tilde v_2=X_3=\frac{t x_2+x_3}{1+t^2},\quad
    \tilde v_3=-X_2=\frac{t x_3-x_2}{1+t^2}.
\end{aligned}
\]

\emph{Spatial acceleration, directly.} Substitute \eqr{2.31} into $\hat a_i$:
$\tilde a_1=e^{t}X_1=x_1$, $\tilde a_2=0$, $\tilde a_3=0$, i.e.\
$\ba=x_1\be_1$.

\emph{Spatial acceleration, via \eqr{2.29} (a check on the material
derivative).} The spatial velocity gradient is
\begin{equation}
   (\tilde v_{i,j})=\begin{pmatrix}
   1&0&0\\[1mm]
   0&\dfrac{t}{1+t^2}&\dfrac{1}{1+t^2}\\[2mm]
   0&-\dfrac{1}{1+t^2}&\dfrac{t}{1+t^2}\end{pmatrix}, \tag{2.33}\label{eq:2.33}
\end{equation}
and the spatial time derivatives $\tilde v_i'=\partial\tilde v_i/\partial t$
(at fixed $\bx$) are $\tilde v_1'=0$ and, by the quotient rule,
\begin{equation}
   \tilde v_2'=\frac{(1-t^2)x_2-2t x_3}{(1+t^2)^2},\qquad
   \tilde v_3'=\frac{(1-t^2)x_3+2t x_2}{(1+t^2)^2}. \tag{2.32}\label{eq:2.32}
\end{equation}
Then $\tilde a_2=\tilde v_2'+\tilde v_{2,j}\tilde v_j
=\tilde v_2'+\tilde v_{2,2}\tilde v_2+\tilde v_{2,3}\tilde v_3$; with
$\tilde v_{2,2}=t/(1+t^2)$, $\tilde v_{2,3}=1/(1+t^2)$,
\[
   \tilde v_{2,2}\tilde v_2+\tilde v_{2,3}\tilde v_3
   =\frac{t(tx_2+x_3)+(tx_3-x_2)}{(1+t^2)^2}
   =\frac{(t^2-1)x_2+2t x_3}{(1+t^2)^2}
   =-\tilde v_2' ,
\]
so $\tilde a_2=0$, matching the direct computation; likewise $\tilde a_1=x_1$,
$\tilde a_3=0$. The two routes agree, which is the point of the example.
\end{example}

\begin{remark}[a correction to the printed $\tilde v_2'$]\label{rem:v2prime}
The material derivative check forces
$\tilde v_2'=[(1-t^2)x_2-2t x_3]/(1+t^2)^2$: only this value makes
$\tilde a_2=\tilde v_2'+\tilde v_{2,j}\tilde v_j$ vanish, as it must since
$\hat a_2=0$. (The text's printed coefficient of $x_2$ in $\tilde v_2'$ is a
misprint; the denominator is squared and the numerator is $1-t^2$, not
$1-2t^2$.)
\end{remark}

%=======================================================================
\section{The deformation gradient}\label{sec:defgrad}
%=======================================================================

Assume the motion $\bchi_\kappa(\bX,t)$ is differentiable in $\bX$. Holding
$t$ fixed, its linear part is a tensor, the \emph{deformation gradient}:
\begin{equation}
   \dl\bx=\bF\,\dl\bX,\qquad \bF=\Grad\bchi_\kappa. \tag{2.34}\label{eq:2.34}
\end{equation}
Allowing $t$ to vary as well,
\begin{equation}
   \dl\bx=\bF\,\dl\bX+\bv\,\dl t. \tag{2.35}\label{eq:2.35}
\end{equation}

\begin{theorem}[The two gradients of a field are related by $\bF$]\label{thm:GradgradF}
For any field $\bu$,
\begin{equation}
   \Grad\bu=(\grad\bu)\bF. \tag{2.37}\label{eq:2.37}
\end{equation}
\end{theorem}

\begin{proof}
Write the increment $\dl\bu$ referentially and spatially, as in \eqr{2.19},
and substitute \eqr{2.35} and the material derivative $\dot\bu=\bu'+(\grad\bu)\bv$:
\begin{equation}
   \dot\bu\,\dl t+(\Grad\bu)\dl\bX
   =[\bu'+(\grad\bu)\bv]\,\dl t+(\grad\bu)\bF\,\dl\bX. \tag{2.36}\label{eq:2.36}
\end{equation}
The $\dl t$ terms cancel by \eqr{2.20}, leaving
$(\Grad\bu)\dl\bX=(\grad\bu)\bF\,\dl\bX$ for every $\dl\bX$; hence
$\Grad\bu=(\grad\bu)\bF$.
\end{proof}

To read off components, differentiate $\bx=\chi_i\be_i$ at fixed $t$ using
$\dl\chi_i=\chi_{i,A}\,\dl X_A=\chi_{i,A}\,\ip{\hbE_A}{\dl\bX}$:
\begin{equation}
   \dl\bx=\chi_{i,A}\be_i\,\ip{\hbE_A}{\dl\bX}
   =(\chi_{i,A}\,\be_i\otimes\hbE_A)\dl\bX, \tag{2.38}\label{eq:2.38}
\end{equation}
so comparison with \eqr{2.34} gives
\begin{equation}
   \bF=F_{iA}\,\be_i\otimes\hbE_A,\qquad F_{iA}=\chi_{i,A}. \tag{2.39}\label{eq:2.39}
\end{equation}
The mixed basis $\{\be_i\otimes\hbE_A\}$ shows $\bF$ is a \emph{two-point
tensor}, a linear map $\hVt\to\Vt$ (exactly the object of Problem~12 of
Chapter~1). Its action on a reference basis vector,
\begin{equation}
   \bF\hbE_B=(F_{iA}\be_i\otimes\hbE_A)\hbE_B
   =F_{iA}(\ip{\hbE_A}{\hbE_B})\be_i=F_{iA}\dd_{AB}\be_i=F_{iB}\be_i, \tag{2.40}\label{eq:2.40}
\end{equation}
gives the component extraction $F_{iA}=\ip{\be_i}{\bF\hbE_A}$ $(2.41)$.

\begin{lemma}[product of two dyads]\label{lem:dyadprod}
For all vectors, $(\ba\otimes\bb)(\bc\otimes\bd)=(\ip\bb\bc)\,\ba\otimes\bd$.
\end{lemma}

\begin{proof}
Act on an arbitrary $\bv$ and unfold using $(\bc\otimes\bd)\bv=(\ip\bd\bv)\bc$
and $(\ba\otimes\bb)\bc=(\ip\bb\bc)\ba$:
\begin{equation}
\begin{aligned}
   [(\ba\otimes\bb)(\bc\otimes\bd)]\bv&=(\ba\otimes\bb)\big[(\ip\bd\bv)\bc\big]
   =(\ip\bd\bv)(\ip\bb\bc)\ba\\
   &=(\ip\bb\bc)(\ip\bd\bv)\ba=(\ip\bb\bc)(\ba\otimes\bd)\bv,
\end{aligned}\tag{2.43--2.44}\label{eq:2.43}
\end{equation}
true for all $\bv$, so the tensors are equal.
\end{proof}

Applying Lemma~\ref{lem:dyadprod} to \eqr{2.37} with $\grad\bu=\tilde u_{i,k}
\be_i\otimes\be_k$ and $\bF=F_{jA}\be_j\otimes\hbE_A$,
\begin{equation}
\begin{aligned}
   \Grad\bu&=(\tilde u_{i,k}\be_i\otimes\be_k)(F_{jA}\be_j\otimes\hbE_A)
   =\tilde u_{i,k}F_{jA}(\ip{\be_k}{\be_j})\be_i\otimes\hbE_A\\
   &=\tilde u_{i,k}\,x_{j,A}\,\dd_{kj}\,\be_i\otimes\hbE_A
   =\tilde u_{i,j}\,x_{j,A}\,\be_i\otimes\hbE_A
   =\hat u_{i,A}\,\be_i\otimes\hbE_A,
\end{aligned}
\tag{2.42}\label{eq:2.42}
\end{equation}
the last step being the ordinary chain rule
$\hat u_{i,A}=\tilde u_{i,j}\,x_{j,A}$. So $\Grad\bu$ is a two-point tensor
too.

\paragraph*{Material curves.}
Let $\bX(S)$, $S$ arclength, parametrise a curve $C\subset\kappa$ of fixed
material points; at time $t^*$ it is convected to
$\bx_*(S)=\bchi_\kappa(\bX(S),t^*)$ (Figure~2.7). The relative position of
two of its points is obtained by integrating the tangent,
\begin{equation}
\begin{aligned}
   \bx_*(S_1)-\bx_*(S_0)&=\be_i\!\int_{S_0}^{S_1}\!x_{*i}'(S)\,\dl S\\
   &=\be_i\!\int_{S_0}^{S_1}\!F_{iA}(X_B(S),t^*)\,X_A'(S)\,\dl S,
\end{aligned}
\tag{2.45--2.46}\label{eq:2.45}
\end{equation}
using $x_{*i}'=\chi_{i,A}X_A'=F_{iA}X_A'$: knowing $\bF$ lets us reconstruct
relative positions of connected material points.

\begin{figure}[htbp]\centering
\begin{tikzpicture}[scale=0.88]
  % curve laid out mostly horizontally; position vector rises steeply: clear angle
  \draw[thick] plot[smooth,tension=0.85] coordinates
    {(-1.5,0.45) (-0.75,1.05) (0.1,0.95) (0.95,1.45) (1.7,1.25)};
  \fill (-1.5,0.45) circle (1.6pt) node[left=2pt] {$S_0$};
  \fill (1.7,1.25) circle (1.6pt) node[right=2pt] {$S_1$};
  \fill (0.1,0.95) circle (1.7pt); \node at (0.12,1.42) {$S$};
  \fill (-0.35,-1.85) circle (1.4pt) node[below=2pt] {$O$};
  \draw[vec] (-0.35,-1.85) -- (0.08,0.86);
  \node at (-0.95,-0.5) {$\bX(S)$};
  \node at (-1.35,1.85) {$C\subset\kappa$};
  \begin{scope}[shift={(6.2,0)}]
  \draw[thick] plot[smooth,tension=0.85] coordinates
    {(-1.5,0.45) (-0.6,0.8) (0.15,1.25) (0.95,1.0) (1.7,1.55)};
  \fill (-1.5,0.45) circle (1.6pt) node[left=2pt] {$S_0$};
  \fill (1.7,1.55) circle (1.6pt) node[right=2pt] {$S_1$};
  \fill (0.15,1.25) circle (1.7pt);
  \fill (-0.35,-1.85) circle (1.4pt) node[below=2pt] {$O$};
  \draw[vec] (-0.35,-1.85) -- (0.12,1.16);
  \node at (-1.05,-0.5) {$\bx_*(S)$};
  \node at (-1.15,1.95) {$C_{t_*}\subset\kappa_{t_*}$};
  \end{scope}
  \draw[mapar] (2.5,1.0) to[bend left=16] (4.05,1.0);
  \node at (3.3,1.65) {$\bchi_\kappa(\cdot,t_*)$};
\end{tikzpicture}
\caption{A material curve and its image in $\kappa_{t_*}$.}\label{fig:2.7}
\end{figure}

\begin{remark}[Figure 2.7]\label{rem:fig7}
Figure~2.7 shows a curve $C$ of fixed material points in $\kappa$ and its
deformed image $C_{t^*}$ in $\kappa_{t^*}$; the arclength $S$ labels the
material points and rides along with them. The tangent transforms by $\bF$,
$\bx_*'=\bF\,\bX'$, the seed of the stretch formula $\mu\bmm=\bF\bM$ of
\S2.6.
\end{remark}

\paragraph*{Volume.}
Invertibility of $\bchi_\kappa$ at each $t$ makes $\bF$ invertible
(inverse function theorem), and invertibility of a tensor is $\det\bF\neq0$,
where, from Chapter~1,
\begin{equation}
   \det\bF=\frac{\tbox{\bF\ba}{\bF\bb}{\bF\bc}}{\tbox\ba\bb\bc} \tag{2.47}\label{eq:2.47}
\end{equation}
for any independent $\{\ba,\bb,\bc\}$ (well defined for the two-point $\bF$
because numerator and denominator each use vectors of a single space).

\begin{remark}[invertibility of $\bF$ is only \emph{local} information]
\label{rem:localinv}
The implication runs one way only. Global invertibility of the motion
$\bchi_\kappa(\cdot,t)$ gives, by the inverse function theorem, an inverse
$\bF^{-1}$ at each $\bX$. But the existence of $\bF^{-1}$ at a
\emph{particular} $\bX$ does \emph{not} imply that $\bchi_\kappa$ is
invertible at all $\bX$: it implies only that $\bchi_\kappa$ is invertible in
some \emph{neighbourhood} of the point considered. Nothing local can prevent
a motion from folding a distant part of the body onto the piece near $\bX$;
$\det\bF\neq0$ everywhere forbids local self-penetration, not global. This is
the same phenomenon as $x\mapsto e^{\mathrm{i}x}$ having nonzero derivative
everywhere while being far from injective, and it is why global
injectivity of the motion is imposed as a separate physical assumption in
Definition~\ref{def:motion} rather than deduced from $\det\bF>0$.
\end{remark}

\noindent Three
material curves through $\bX$ carry tangents $\dl\bX,\dl\bY,\dl\bZ$ spanning
a reference parallelepiped of volume $\dl V=\tbox{\dl\bX}{\dl\bY}{\dl\bZ}$
(Figure~2.8); their images $\dl\bx=\bF\dl\bX$ etc.\ span
$\dl v=\tbox{\dl\bx}{\dl\by}{\dl\bz}$. Taking $\{\ba,\bb,\bc\}=\{\dl\bX,\dl\bY,\dl\bZ\}$
in \eqr{2.47},
\begin{equation}
\begin{aligned}
   \dl v&=\tbox{\bF\dl\bX}{\bF\dl\bY}{\bF\dl\bZ}
        =\det\bF\;\tbox{\dl\bX}{\dl\bY}{\dl\bZ}\\
   &=J\,\dl V,\qquad J=\det\bF>0.
\end{aligned}
\tag{2.48--2.50}\label{eq:2.48}
\end{equation}

\begin{figure}[htbp]\centering
\begin{tikzpicture}[scale=0.92]
% draw a full parallelepiped from an origin and three edge vectors
\newcommand{\ppiped}[6]{% #1 origin  #2,#3,#4 edges (u,v,w)  #5 edge style #6 arrow style
  \begin{scope}[shift={#1}]
    \coordinate (o) at (0,0); \coordinate (U) at #2;
    \coordinate (V) at #3;    \coordinate (W) at #4;
    \draw[#5] (o)--(U) (o)--(V) (o)--(W);
    \draw[#5] (U)--($(U)+(V)$) (U)--($(U)+(W)$);
    \draw[#5] (V)--($(U)+(V)$) (V)--($(V)+(W)$);
    \draw[#5] (W)--($(U)+(W)$) (W)--($(V)+(W)$);
    \draw[#5] ($(U)+(V)$)--($(U)+(V)+(W)$)
              ($(U)+(W)$)--($(U)+(V)+(W)$)
              ($(V)+(W)$)--($(U)+(V)+(W)$);
  \end{scope}}
  % reference
  \ppiped{(0,0)}{(1.45,-0.3)}{(0.35,1.25)}{(0.95,0.42)}{gray!65,thin}{}
  \draw[vec] (0,0) -- (1.45,-0.3) node[below=1pt] {$\dl\bX$};
  \draw[vec] (0,0) -- (0.35,1.25) node[above left=-1pt] {$\dl\bY$};
  \draw[vec] (0,0) -- (0.95,0.42) node[right=1pt] {$\dl\bZ$};
  \node at (1.9,1.75) {$\dl V$};  \node at (-1.15,1.35) {$\kappa$};
  % current (sheared and stretched)
  \begin{scope}[shift={(5.5,-0.1)}]
  \ppiped{(0,0)}{(1.75,0.15)}{(0.05,1.35)}{(1.05,0.65)}{gray!65,thin}{}
  \draw[vec] (0,0) -- (1.75,0.15) node[below=1pt] {$\dl\bx$};
  \draw[vec] (0,0) -- (0.05,1.35) node[above left=-1pt] {$\dl\by$};
  \draw[vec] (0,0) -- (1.05,0.65) node[right=1pt] {$\dl\bz$};
  \node at (2.35,1.95) {$\dl v$};  \node at (3.35,1.35) {$\kappa_t$};
  \end{scope}
  \draw[mapar] (2.9,1.95) to[bend left=16] (4.85,1.95);
  \node at (3.9,2.55) {$\bF$};
  \node at (3.9,-1.05) {$\dl v=J\,\dl V,\qquad J=\det\bF$};
\end{tikzpicture}
\caption{Three material curves intersecting at a material point.}\label{fig:2.8}
\end{figure}

\begin{remark}[Figure 2.8, and $J$ as a volume ratio]\label{rem:fig8}
Figure~2.8 shows the reference parallelepiped (left, $\kappa$) and its image
(right, $\kappa_t$): $J=\det\bF$ is the factor by which local volume is
scaled, exactly the ``oriented volume ratio'' reading of the determinant
from Chapter~1. If the reference configuration is \emph{occupiable} (could
be physically occupied), then $J>0$; if not, one only knows $J=|\det\bF|$.
Incompressibility is the constraint $J\equiv1$.
\end{remark}

\paragraph*{Area (the Piola--Nanson formula).}
Let $\dl\bX,\dl\bY$ span the tangent plane of a material surface at $\bX$,
of area $\dl A=|\dl\bX\times\dl\bY|$, with unit normal $\bN$ defined by
\begin{equation}
   \bN\,\dl A=\dl\bX\times\dl\bY \tag{2.51}\label{eq:2.51}
\end{equation}
(Figure~2.9). The deformed edges $\dl\bx=\bF\dl\bX$, $\dl\by=\bF\dl\bY$ span
a parallelogram of area $\dl a$ with unit normal $\bn$, $\bn\,\dl a=\dl\bx
\times\dl\by$. Using the cofactor $\bF^{*}$, defined (as in Chapter~1) by
$\bF^{*}(\ba\times\bb)=\bF\ba\times\bF\bb$,
\begin{equation}
   \bn\,\dl a=\bF\dl\bX\times\bF\dl\bY=\bF^{*}(\dl\bX\times\dl\bY)
   =\bF^{*}\bN\,\dl A, \tag{2.52--2.54}\label{eq:2.52}
\end{equation}
the \emph{Piola--Nanson formula}, in components $n_i\,\dl a=F^{*}_{iB}N_B\,\dl A$.
By Problem~22(c) of Chapter~1, $\bA^{t}\bA^{*}=(\det\bA)\I$, valid also for
the two-point $\bF$, so
\begin{equation}
   \bF^{*}=J\bF^{-t}. \tag{2.55}\label{eq:2.55}
\end{equation}

\begin{figure}[htbp]\centering
\begin{tikzpicture}[scale=0.92]
  % reference patch drawn in perspective: edges nearly horizontal, normal straight up
  \fill[gray!12] (0,0)--++(1.55,-0.12)--++(0.62,0.55)--++(-1.55,0.12)--cycle;
  \draw[gray!70,thin] (0,0)--++(1.55,-0.12)--++(0.62,0.55)--++(-1.55,0.12)--cycle;
  \draw[vec] (0,0) -- ++(1.55,-0.12) node[below right=-1pt and -2pt] {$\dl\bX$};
  \draw[vec] (0,0) -- ++(0.62,0.55) node[left=1pt] {$\dl\bY$};
  \draw[vec] (1.09,0.22) -- ++(0,1.25) node[above=1pt] {$\bN$};
  \draw[gray!70,thin] (1.09,0.22) ++(-0.16,0) -- ++(0.16,0.16) -- ++(0.16,-0.16);
  \node at (1.62,0.34) {$\dl A$};  \node at (-0.85,1.15) {$\kappa$};
  % current patch
  \begin{scope}[shift={(5.4,0)}]
  \fill[gray!12] (0,0)--++(1.9,0.22)--++(0.35,0.72)--++(-1.9,-0.22)--cycle;
  \draw[gray!70,thin] (0,0)--++(1.9,0.22)--++(0.35,0.72)--++(-1.9,-0.22)--cycle;
  \draw[vec] (0,0) -- ++(1.9,0.22) node[below right=-1pt and -2pt] {$\dl\bx$};
  \draw[vec] (0,0) -- ++(0.35,0.72) node[left=1pt] {$\dl\by$};
  \draw[vec] (1.12,0.47) -- ++(0.28,1.25) node[above=1pt] {$\bn$};
  \node at (1.95,0.62) {$\dl a$};  \node at (2.9,1.15) {$\kappa_t$};
  \end{scope}
  \draw[mapar] (2.85,1.55) to[bend left=15] (4.9,1.6);
  \node at (3.85,2.2) {$\bF$};
  \node at (3.9,-1.15) {$\bn\,\dl a=\bF^{*}\bN\,\dl A=J\bF^{-t}\bN\,\dl A$};
  % inset: area of a parallelogram
  \begin{scope}[shift={(2.15,-3.35)},scale=0.95]
  \fill[gray!12] (0,0)--++(1.5,0)--++(0.65,0.85)--++(-1.5,0)--cycle;
  \draw[gray!70,thin] (0,0)--++(1.5,0)--++(0.65,0.85)--++(-1.5,0)--cycle;
  \draw[vec] (0,0) -- (1.5,0) node[below=1pt] {$\ba$};
  \draw[vec] (0,0) -- (0.65,0.85) node[above left=-2pt] {$\bb$};
  \draw (0.45,0) arc (0:52.6:0.45); \node at (0.72,0.22) {\small$\alpha$};
  \node at (1.05,-0.62) {\small area $=|\ba||\bb|\sin\alpha=|\ba\times\bb|$};
  \end{scope}
\end{tikzpicture}
\caption{Deformation of oriented surface measure.}\label{fig:2.9}
\end{figure}

\begin{remark}[Figure 2.9: area is not a passive rider]\label{rem:fig9}
Figure~2.9 shows the reference patch with edges $\dl\bX,\dl\bY$, area $\dl A$,
normal $\bN$, and its deformed image with $\dl\bx,\dl\by$, area $\dl a$,
normal $\bn$; the small inset recalls $\text{area}=|\ba||\bb|\sin\alpha
=|\ba\times\bb|$ from Chapter~1. The point of $\bn\,\dl a=J\bF^{-t}\bN\,\dl A$
is that the normal does \emph{not} transform by $\bF$ (that would be the
tangent's rule): it transforms by the cofactor $\bF^{*}=J\bF^{-t}$, because a
normal is defined by a cross product, and cross products go with cofactors.
This single formula converts every surface traction and flux between
reference and current descriptions in the balance laws.
\end{remark}

%=======================================================================
\section{Mass and density}\label{sec:mass}
%=======================================================================

Consider a sub-body $S\subset B$, occupying a part $P=\bkappa(S)\subset\kappa$
of the reference configuration and a part $P_t=\bchi_\kappa(P,t)\subset\kappa_t$
of the current one. Mass is assumed absolutely continuous in volume (more
volume, more mass) and strictly positive on regions of positive volume, so
there are positive \emph{density} fields $\rho(\bx,t)$ and $\rho_\kappa(\bX,t)$
with
\begin{equation}
   m(S,t)=\int_{P_t}\rho(\bx,t)\,\dl v=\int_{P}\rho_\kappa(\bX,t)\,\dl V. \tag{2.56}\label{eq:2.56}
\end{equation}
(Conservation of mass, treated later, will make $m$ independent of $t$.)

Change variables in the first integral by $\bx=\bchi_\kappa(\bX,t)$,
$\dl v=J\,\dl V$ from \eqr{2.50}:
\begin{equation}
   \int_{P_t}\rho(\bx,t)\,\dl v=\int_{P}\rho(\bchi_\kappa(\bX,t),t)\,J(\bX,t)\,\dl V,
   \tag{2.57}\label{eq:2.57}
\end{equation}
so equality of the two forms in \eqr{2.56} becomes
\begin{equation}
   \int_{P}\big[\rho_\kappa(\bX,t)-\rho(\bchi_\kappa(\bX,t),t)\,J(\bX,t)\big]\,\dl V=0
   \qquad\text{for all }P\subset\kappa. \tag{2.58}\label{eq:2.58}
\end{equation}
It is \emph{sufficient} for this that the bracket vanish pointwise,
\begin{equation}
   \rho_\kappa=\rho J\quad\text{at every }\bX\in\kappa. \tag{2.59}\label{eq:2.59}
\end{equation}
That it is also \emph{necessary}, when the integrand is continuous, is the
content of the localization theorem.

\begin{theorem}[Localization]\label{thm:localization}
Let $\varphi(\bX)$ be continuous and suppose
\begin{equation}
   \int_{P}\varphi(\bX)\,\dl V=0\qquad\text{for all }P\subset\kappa. \tag{2.60}\label{eq:2.60}
\end{equation}
Then
\begin{equation}
   \varphi(\bX)=0\quad\text{at every }\bX\in\kappa. \tag{2.61}\label{eq:2.61}
\end{equation}
\textup{(The converse statement is obvious: if $\varphi$ vanishes
identically its integral over every $P$ is zero.)}
\end{theorem}

\begin{proof}
Fix $\bX_0\in\kappa$ and take $P=R_\eps$, the ball of radius $\eps$ and
volume $V_\eps$ centred at $\bX_0$ and contained in $\kappa$. Consider the
average discrepancy
\begin{equation}
   I_\eps=\bigg|\frac{1}{V_\eps}\int_{R_\eps}[\varphi(\bX_0)-\varphi(\bX)]\,\dl V\bigg|.
   \tag{2.62}\label{eq:2.62}
\end{equation}
By the triangle inequality for integrals, then bounding the integrand by its
maximum over the (compact) ball,
\begin{equation}
\begin{aligned}
   I_\eps&\le\frac{1}{V_\eps}\int_{R_\eps}|\varphi(\bX_0)-\varphi(\bX)|\,\dl V\\
   &\le\frac{1}{V_\eps}\int_{R_\eps}\max_{\bX\in R_\eps}
       |\varphi(\bX_0)-\varphi(\bX)|\,\dl V\\
   &=\max_{\bX\in R_\eps}|\varphi(\bX_0)-\varphi(\bX)|,
\end{aligned}
\tag{2.63}\label{eq:2.63}
\end{equation}
the last step because the maximum is a constant and $\int_{R_\eps}\dl V=V_\eps$.
A continuous function on a compact set attains its maximum, and by continuity
of $\varphi$ that maximum tends to $0$ as the ball shrinks:
\begin{equation}
   0\le I_\eps\le\max_{\bX\in R_\eps}|\varphi(\bX_0)-\varphi(\bX)|\ \longrightarrow\ 0
   \quad(\eps\to0). \tag{2.64}\label{eq:2.64}
\end{equation}
But the hypothesis \eqr{2.60} makes $\int_{R_\eps}\varphi(\bX)\,\dl V=0$, so
\begin{equation}
   \varphi(\bX_0)=\frac{1}{V_\eps}\int_{R_\eps}\varphi(\bX_0)\,\dl V
   =\frac{1}{V_\eps}\int_{R_\eps}[\varphi(\bX_0)-\varphi(\bX)]\,\dl V
   \xrightarrow{\ \eps\to0\ }0, \tag{2.65}\label{eq:2.65}
\end{equation}
i.e.\ $|\varphi(\bX_0)|=\lim_{\eps\to0}I_\eps=0$. Since $\bX_0$ was
arbitrary, $\varphi\equiv0$.
\end{proof}

Applying the theorem to the continuous bracket in \eqr{2.58} yields the
pointwise mass relation \eqr{2.59}, $\rho_\kappa=\rho J$. Continuity of that
bracket is guaranteed by the assumed differentiability of $\bchi_\kappa$
(which makes $\bF$, and hence $J=\det\bF$, continuous) together with
continuity of $\rho$.

\begin{remark}[what localization really says]\label{rem:localuse}
The theorem is the bridge from \emph{integral} (global) balance laws to
\emph{pointwise} (local) field equations, and it is used again and again:
conservation of mass, momentum, and energy are all first stated as ``the
integral over every sub-body vanishes'', and localization converts each into
a differential equation holding at every point. Intuitively, if an average
over \emph{every} shrinking ball around $\bX_0$ is forced to be zero, the
value at $\bX_0$ cannot be anything but zero, since a continuous function is
nearly constant on a small enough ball and its average there nearly equals
its central value. The relation $\rho_\kappa=\rho J$ itself is just
conservation of mass in disguise: the same matter in a smaller current
volume ($J<1$) must be denser, $\rho=\rho_\kappa/J$.
\end{remark}

%=======================================================================
\section{Material curves and the shifter}\label{sec:shifter}
%=======================================================================

Return to the material curve $C\subset\kappa$ and its convected image
$C_{t_*}=\bchi_\kappa(C,t_*)\subset\kappa_{t_*}$ (Figure~2.10). The material
point at arclength station $S$ on $C$ occupies $\bx_*(S)=\bchi_\kappa(\bX(S),t_*)$.
Differentiating in $S$ at fixed $t_*$, by the chain rule and \eqr{2.34},
\begin{equation}
   \bx_*'(S)=\bF(\bX(S),t_*)\,\bM(S),\qquad \bM(S)=\bX'(S), \tag{2.66}\label{eq:2.66}
\end{equation}
$\bM$ being a \emph{unit} tangent to $C$ (because $S$ is arclength,
$|\bX'|=1$). The image $\bx_*'$ is tangent to $C_{t_*}$ but need not be a
unit vector; writing $\bx_*'=\mu\bmm$ with $|\bmm|=1$ and $\mu=|\bx_*'|$ the
local \emph{stretch},
\begin{equation}
   \mu\bmm=\bF\bM. \tag{2.67}\label{eq:2.67}
\end{equation}

\begin{figure}[htbp]\centering
\begin{tikzpicture}[scale=0.9]
  % curve thin+gray so the tangent stands out
  \draw[gray!75,thick] plot[smooth,tension=0.9] coordinates
    {(-1.3,-1.35) (-0.75,-0.35) (-0.05,0.25) (0.5,1.15) (1.05,1.95)};
  \fill (-0.05,0.25) circle (1.9pt); \node at (-0.5,0.42) {$S$};
  \draw[vec,black] (-0.05,0.25) -- ++(0.55,0.9);
  \node at (0.88,0.82) {$\bM$};
  \node[gray!75] at (-1.35,1.3) {$C\subset\kappa$};
  \node at (0.35,-0.55) {\small$|\bM|=1$};
  \begin{scope}[shift={(5.6,0)}]
  \draw[gray!75,thick] plot[smooth,tension=0.9] coordinates
    {(-1.3,-1.35) (-0.4,-0.9) (0.05,0.1) (0.3,1.05) (1.0,1.95)};
  \fill (0.05,0.1) circle (1.9pt);
  \draw[vec,black] (0.05,0.1) -- ++(0.28,1.3);
  \node at (0.95,1.42) {$\mu\bmm$};
  \node[gray!75] at (-1.15,-1.35) {$C_{t_*}\subset\kappa_{t_*}$};
  \node at (0.35,-0.65) {\small$|\bmm|=1,\ \mu=|\bF\bM|$};
  \end{scope}
  \draw[mapar] (1.9,1.5) to[bend left=16] (3.6,1.5);
  \node at (2.75,2.1) {$\mu\bmm=\bF\bM$};
\end{tikzpicture}
\caption{Orientation of material curves at point $p$ with arclength $S$ in $\kappa$.}
\label{fig:2.10}
\end{figure}

\begin{remark}[Figure 2.10: what $\bF$ does to a fibre]\label{rem:fig10}
Figure~2.10 shows the curve $C\subset\kappa$ with station $S$, position
$\bX(S)$ and unit tangent $\bM$, beside its image $C_{t_*}\subset\kappa_{t_*}$
with $\bx_*(S)=\bchi_\kappa(\bX(S),t_*)$. Equation \eqr{2.67} is the sentence
``$\bF$ takes the unit material fibre $\bM$ and returns the deformed fibre,
whose length $\mu$ is the stretch and whose direction $\bmm$ is the new
orientation''. Every deformation measure in \S\ref{sec:cauchygreen} is
extracted from this one statement by taking lengths and angles of $\bF\bM$.
\end{remark}

\medskip
Interpreting \eqr{2.67} for an \emph{undeformed} body raises a genuine
conceptual difficulty. We stipulated $\bX\in\hVt$ but $\kappa_t\subset\Et$
for every $t$, including $t_0$; so if we want to say ``at $t_0$ the body is
where its reference copy is'', we are comparing vectors from two different
spaces, which is meaningless. The repair is a fixed two-point tensor that
carries one space onto the other.

\begin{definition}[The shifter]\label{def:shifter}
Introduce fixed orthonormal vectors $\bE_A\in\Vt$ and the uniform two-point
tensor
\begin{equation}
   \vek{1}=\bE_A\otimes\hbE_A . \tag{2.68}\label{eq:2.68}
\end{equation}
If $\{\hbE_A\}$ and $\{\bE_A\}$ have the same handedness, $\vek{1}$
rotates the former onto the latter and is a two-point rotation; otherwise it
is merely orthogonal. Resolving $\bE_A\in\Vt$ on $\{\be_i\}$ by the
orthonormal expansion of Chapter~1,
\begin{equation}
   \bE_A=\dd_{iA}\be_i,\qquad \dd_{iA}=\ip{\be_i}{\bE_A}, \tag{2.69}\label{eq:2.69}
\end{equation}
(the dot product is legitimate: both factors lie in $\Vt$), giving the
standard two-point decomposition
\begin{equation}
   \vek{1}=\dd_{iA}\,\be_i\otimes\hbE_A . \tag{2.70}\label{eq:2.70}
\end{equation}
The $\dd_{iA}$ are Kronecker deltas if and only if $\{\bE_A\}=\{\be_i\}$;
otherwise they are the cosines of the fixed angles between the two sets.
\end{definition}

\begin{proposition}[Algebra of the shifter]\label{prop:shifter}
With $\I$ and $\hat\I$ the identities of $\Vt$ and $\hVt$,
\begin{equation}
   \vek{1}=\vek{1}\hat\I=\I\vek{1},\qquad
   \vek{1}^{t}\vek{1}=\hat\I,\qquad
   \vek{1}\vek{1}^{t}=\I, \tag{2.71--2.73}\label{eq:2.71}
\end{equation}
equivalently, in components,
\begin{equation}
   \dd_{iA}\dd_{iB}=\dd_{AB},\qquad \dd_{iA}\dd_{jA}=\dd_{ij}. \tag{2.74}\label{eq:2.74}
\end{equation}
\end{proposition}

\begin{proof}
For \eqr{2.71}: $\bE_A=\I\bE_A$ and $\hbE_A=\hat\I\hbE_A$, so by the rules
$\bA(\ba\otimes\bb)=(\bA\ba)\otimes\bb$ and
$(\ba\otimes\bb)\bA=\ba\otimes(\bA^{t}\bb)$ (Problems 17(b),(c) of
Chapter~1),
\[
\begin{aligned}
   \I\vek{1}&=\I(\bE_A\otimes\hbE_A)=(\I\bE_A)\otimes\hbE_A=\vek{1},\\
   \vek{1}\hat\I&=\bE_A\otimes(\hat\I^{t}\hbE_A)=\bE_A\otimes\hbE_A=\vek{1}.
\end{aligned}
\]
For \eqr{2.72}, transpose \eqr{2.68} by Problem 17(a) of Chapter~1, which gives
$(\ba\otimes\bb)^{t}=\bb\otimes\ba$, and then multiply using
Lemma~\ref{lem:dyadprod}:
\begin{equation}
\begin{aligned}
   \vek{1}^{t}\vek{1}&=(\hbE_A\otimes\bE_A)(\bE_B\otimes\hbE_B)
   =(\ip{\bE_A}{\bE_B})\,\hbE_A\otimes\hbE_B\\
   &=\dd_{AB}\,\hbE_A\otimes\hbE_B=\hbE_A\otimes\hbE_A=\hat\I .
\end{aligned}\tag{2.72}
\end{equation}
Identically, $\vek{1}\vek{1}^{t}=(\ip{\hbE_A}{\hbE_B})\bE_A\otimes\bE_B
=\bE_A\otimes\bE_A=\I$. In components, \eqr{2.72} reads
$\dd_{iA}\dd_{iB}=\dd_{AB}$ and \eqr{2.73} reads $\dd_{iA}\dd_{jA}=\dd_{ij}$,
which is \eqr{2.74}.
\end{proof}

Acting with the shifter on \eqr{2.23} produces a \emph{spatial} copy of the
reference position,
\begin{equation}
   \vek{1}\bX=X_A\,\vek{1}\hbE_A=X_A\bE_A\in\Vt, \tag{2.75}\label{eq:2.75}
\end{equation}
and the statement $\kappa=\kappa_{t_0}$ is then properly written
\begin{equation}
   \bchi_\kappa(\bX,t_0)=\vek{1}\bX \tag{2.76}\label{eq:2.76}
\end{equation}
(taking the origins of Figure~2.6 to coincide, without loss of generality).
This makes the \emph{displacement} well defined:
\begin{equation}
   \bu(\bX,t)=\bchi_\kappa(\bX,t)-\vek{1}\bX\ \in\Vt,\qquad
   \hat u_i(X_B,t)=\chi_i(X_B,t)-\dd_{iA}X_A. \tag{2.77--2.78}\label{eq:2.77}
\end{equation}

\begin{remark}[why the shifter is not pedantry]\label{rem:shifteruse}
The definition of displacement found in most of the literature omits
$\vek{1}$ and writes $\bu=\bchi_\kappa-\bX$. That expression subtracts a
vector of $\Vt$ from a vector of $\hVt$: an impermissible operation between
different vector spaces, and a documented source of confusion in using
displacement to characterise deformation. The shifter is the ``two-point
identity tensor'': when $\bE_A$ and $\hbE_A$ are aligned it simply
\emph{parallel transports} a vector from one copy of space to the other and
changes nothing else, so its components $\dd_{iA}$ then really are Kronecker
deltas and every classical formula is recovered. Its only job is to make the
bookkeeping honest, exactly as the caret did for functions
(Remark~\ref{rem:decorations}). The shifter and the displacement are put to
substantive use later, in the book's Chapters~7 and~12.
\end{remark}

\begin{proposition}[Undeformed state]\label{prop:undeformed}
If \eqr{2.76} holds then
\begin{equation}
   \bF(\bX,t_0)=\vek{1}, \qquad F_{iA}=\dd_{iA}, \tag{2.79}\label{eq:2.79}
\end{equation}
and consequently $\mu=1$ for every material curve:
\begin{equation}
   \mu\bmm=\vek{1}\bM,\qquad
   \mu=\sqrt{\ip{\vek{1}\bM}{\vek{1}\bM}}=\sqrt{\ip\bM\bM}=1. \tag{2.80--2.81}\label{eq:2.80}
\end{equation}
Moreover, in the general case \eqr{2.67} forces $\mu>0$.
\end{proposition}

\begin{proof}
\eqr{2.79}: $\bF=\Grad\bchi_\kappa$ and $\bchi_\kappa(\bX,t_0)=\vek{1}\bX$ is
\emph{linear} in $\bX$ with the uniform (position independent) coefficient
$\vek{1}$, so its gradient is $\vek{1}$ itself; in components
$F_{iA}=\partial(\dd_{jB}X_B)/\partial X_A\,\big|_{i=j}=\dd_{iA}$.
Substituting into \eqr{2.67} gives \eqr{2.80}, and since the shifter is
orthogonal, $\ip{\vek{1}\bM}{\vek{1}\bM}
=\ip{\bM}{\vek{1}^{t}\vek{1}\bM}=\ip\bM{\hat\I\bM}=\ip\bM\bM=1$ by
\eqr{2.72}, so $\mu=1$: an undeformed body stretches no fibre.
Finally $\mu>0$ always, for $\mu=0$ would give $\bF\bM=\bze$ with
$\bM\neq\bze$ (as $|\bM|=1$), contradicting invertibility of $\bF$.
\end{proof}

%=======================================================================
\section{The Cauchy--Green deformation tensors and the Lagrange--Euler
strain tensors}\label{sec:cauchygreen}
%=======================================================================

The stretch is the positive square root of $\mu^2$, and \eqr{2.67} turns that
into a quadratic form in $\bM$. Using only the definitions of the transpose
and of the product of two tensors (Chapter~1),
\begin{equation}
   \mu^{2}=|\mu\bmm|^{2}=\ip{\bF\bM}{\bF\bM}
   =\ip\bM{\bF^{t}(\bF\bM)}=\ip\bM{(\bF^{t}\bF)\bM}, \tag{2.82}\label{eq:2.82}
\end{equation}
so that
\begin{equation}
   0<\mu^{2}=\ip\bM{\bC\bM},\qquad \bC=\bF^{t}\bF \tag{2.83}\label{eq:2.83}
\end{equation}
defines the \emph{right Cauchy--Green deformation tensor}.

\begin{proposition}[Properties of $\bC$]\label{prop:C}
$\bC$ is a \emph{referential} tensor ($\hVt\to\hVt$), symmetric, and
positive definite; its components are
\begin{equation}
\begin{aligned}
   &\bC=C_{AB}\,\hbE_A\otimes\hbE_B,\qquad
   C_{AB}=F_{iA}F_{iB}=C_{BA},\\
   &\mu^{2}=C_{AB}M_AM_B.
\end{aligned}
\tag{2.86--2.88}\label{eq:2.86}
\end{equation}
\end{proposition}

\begin{proof}
\emph{Domain.} $\bF:\hVt\to\Vt$, hence $\bF^{t}:\Vt\to\hVt$, hence
$\bC=\bF^{t}\bF$ maps $\hVt$ into itself.
\emph{Positive definiteness.} For any $\ba\neq\bze$ in $\hVt$, moving $\bF^t$
across the dot product by the definition of the transpose,
\begin{equation}
   \ip\ba{\bC\ba}=\ip\ba{\bF^{t}\bF\ba}=\ip{\bF\ba}{\bF\ba}=|\bF\ba|^{2}>0,
   \tag{2.84}\label{eq:2.84}
\end{equation}
strictly, because $\bF$ is invertible ($\det\bF\neq0$) so $\bF\ba\neq\bze$.
\emph{Symmetry,} using $(\bB\bA)^{t}=\bA^{t}\bB^{t}$ (Problem 20 of
Chapter~1) and $(\bF^{t})^{t}=\bF$:
\begin{equation}
   \bC^{t}=(\bF^{t}\bF)^{t}=\bF^{t}(\bF^{t})^{t}=\bF^{t}\bF=\bC. \tag{2.85}\label{eq:2.85}
\end{equation}
\emph{Components.} With $\bF=F_{iA}\be_i\otimes\hbE_A$ and
$\bF^{t}=F_{iA}\hbE_A\otimes\be_i$ (Problem 17(a) of Chapter~1), multiply by
Lemma~\ref{lem:dyadprod}:
\begin{equation}
\begin{aligned}
   \bC&=(F_{iA}\hbE_A\otimes\be_i)(F_{jB}\be_j\otimes\hbE_B)
   =F_{iA}F_{jB}(\ip{\be_i}{\be_j})\,\hbE_A\otimes\hbE_B\\
   &=F_{iA}F_{jB}\dd_{ij}\,\hbE_A\otimes\hbE_B
   =C_{AB}\,\hbE_A\otimes\hbE_B,\qquad C_{AB}=F_{iA}F_{iB},
\end{aligned}\tag{2.86--2.87}
\end{equation}
manifestly symmetric since $F_{iA}F_{iB}=F_{iB}F_{iA}$. Finally
\begin{equation}
   \mu^{2}=\ip\bM{(C_{AB}\hbE_A\otimes\hbE_B)\bM}
   =\ip\bM{C_{AB}M_B\hbE_A}=C_{AB}M_AM_B, \tag{2.88}
\end{equation}
so $(C_{AB})$ is a positive definite matrix.
\end{proof}

Inverting \eqr{2.67} gives the companion, spatial, measure. Applying $\bF^{-1}$
to $\mu\bmm=\bF\bM$,
\begin{equation}
\begin{aligned}
   \mu\bF^{-1}\bmm&=\bF^{-1}(\bF\bM)=(\bF^{-1}\bF)\bM=\hat\I\bM=\bM,\\
   \text{i.e.}\quad \mu^{-1}\bM&=\bF^{-1}\bmm,
\end{aligned}
\tag{2.89--2.90}\label{eq:2.89}
\end{equation}
and taking the squared length of the left side (recall $|\bM|=1$),
\begin{equation}
\begin{aligned}
   \mu^{-2}&=\big|\mu^{-1}\bM\big|^{2}=\ip{\bF^{-1}\bmm}{\bF^{-1}\bmm}\\
   &=\ip\bmm{(\bF^{-1})^{t}\bF^{-1}\bmm}
   =\ip\bmm{\bF^{-t}\bF^{-1}\bmm}=\ip\bmm{\bB^{-1}\bmm},
\end{aligned}
\tag{2.91}\label{eq:2.91}
\end{equation}
where
\begin{equation}
   \bB=\bF\bF^{t} \tag{2.92}\label{eq:2.92}
\end{equation}
is the \emph{left Cauchy--Green deformation tensor}; indeed
$\bB^{-1}=(\bF\bF^{t})^{-1}=\bF^{-t}\bF^{-1}$ by Problems 18--19 of
Chapter~1.

\begin{proposition}[Properties of $\bB$]\label{prop:B}
$\bB$ is a \emph{spatial} tensor ($\Vt\to\Vt$), symmetric, with
\begin{equation}
   \bB=B_{ij}\,\be_i\otimes\be_j,\qquad B_{ij}=F_{iA}F_{jA}=B_{ji}, \tag{2.93--2.94}\label{eq:2.93}
\end{equation}
and $\bB^{-1}$ is positive definite.
\end{proposition}

\begin{proof}
$\bF^{t}:\Vt\to\hVt$ then $\bF:\hVt\to\Vt$, so $\bB$ maps $\Vt$ to itself;
symmetry is $\bB^{t}=(\bF\bF^{t})^{t}=\bF\bF^{t}=\bB$. Components, again by
Lemma~\ref{lem:dyadprod}:
\begin{equation}
\begin{aligned}
   \bB&=(F_{iA}\be_i\otimes\hbE_A)(F_{jB}\hbE_B\otimes\be_j)
   =F_{iA}F_{jB}(\ip{\hbE_A}{\hbE_B})\,\be_i\otimes\be_j\\
   &=F_{iA}F_{jB}\dd_{AB}\,\be_i\otimes\be_j=B_{ij}\be_i\otimes\be_j,
   \qquad B_{ij}=F_{iA}F_{jA}.
\end{aligned}\tag{2.93--2.94}
\end{equation}
Positive definiteness of $\bB^{-1}$: for $\ba\neq\bze$ in $\Vt$,
\begin{equation}
   \ip\ba{\bB^{-1}\ba}=\ip\ba{\bF^{-t}\bF^{-1}\ba}
   =\ip{\bF^{-1}\ba}{\bF^{-1}\ba}=\big|\bF^{-1}\ba\big|^{2}>0, \tag{2.95}\label{eq:2.95}
\end{equation}
since $\bF^{-1}$ is invertible (with inverse $\bF$). In Chapter~3 this is
shown equivalent to positive definiteness of $\bB$ itself.
\end{proof}

\begin{remark}[$\bC$ and $\bB$: which space you measure in]\label{rem:CBintuition}
$\bF$ has one foot in each configuration, so it cannot be symmetric and
carries rotation as well as stretch. Pairing $\bF$ with its transpose
cancels the rotation and lands the result wholly in one space:
$\bC=\bF^{t}\bF$ lives in the \emph{reference} space and is fed reference
fibres $\bM$, while $\bB=\bF\bF^{t}$ lives in the \emph{current} space and
is fed current fibres $\bmm$. Concretely, \eqr{2.83} says ``give me the
undeformed direction and I return the square of the stretch'', whereas
\eqr{2.91} says ``give me the deformed direction and I return the inverse
square of the stretch''. If $\bF$ is a pure rotation, $\bF^{t}\bF=\hat\I$
and no fibre stretches, which is exactly why these products, and not $\bF$,
are the measures of \emph{distortion}.
\end{remark}

\begin{figure}[htbp]\centering
\begin{tikzpicture}[scale=1.05]
  % ---- reference: unit circle of directions M
  \draw[gray!55,thin] (-1.35,0)--(1.35,0); \draw[gray!55,thin] (0,-1.35)--(0,1.35);
  \draw[thick] (0,0) circle (1);
  \foreach \a in {15,45,...,345}
     {\draw[gray!55,-{Stealth[length=1.4mm]}] (0,0)--(\a:1);}
  \draw[vec,black] (0,0)--(0:1);  \node at (1.32,0.28) {$\bM_{\!1}$};
  \draw[vec,black] (0,0)--(90:1); \node at (0.36,1.24) {$\bM_{\!2}$};
  \node at (0,-1.85) {\small unit directions $|\bM|=1$ in $\kappa$};
  % ---- map arrow
  \draw[mapar] (1.9,0.4) to[bend left=14] (3.3,0.4);
  \node at (2.6,1.0) {$\bF$};
  % ---- current: ellipse of images FM
  \begin{scope}[shift={(5.6,0)},rotate=22]
  \draw[gray!55,thin] (-1.85,0)--(1.85,0); \draw[gray!55,thin] (0,-0.95)--(0,0.95);
  \draw[thick] (0,0) ellipse (1.62 and 0.7);
  \foreach \a in {15,45,...,345}
     {\draw[gray!55,-{Stealth[length=1.4mm]}] (0,0)--({1.62*cos(\a)},{0.7*sin(\a)});}
  \draw[vec,black] (0,0)--(1.62,0);  \node at (1.5,-0.42) {$\mu_1\bmm_1$};
  \draw[vec,black] (0,0)--(0,0.7);   \node at (-0.78,0.95) {$\mu_2\bmm_2$};
  \end{scope}
  \node at (5.6,-1.85) {\small images $\bF\bM$: an ellipse};
\end{tikzpicture}
\\[1.8ex]
\parbox{0.93\textwidth}{\footnotesize\textsc{Supplementary Figure B.} What
$\bC=\bF^{t}\bF$ measures. Feed $\bF$ every unit direction $\bM$ at a
material point: the images $\bF\bM$ sweep out an \emph{ellipse} (an
ellipsoid in three dimensions), because $\mu^{2}=\ip\bM{\bC\bM}$ is a
quadratic form in $\bM$. The stretch $\mu$ is the length of the image, so it
depends on direction; its largest and smallest values occur along the
principal axes of the ellipse, which are the eigenvectors of $\bC$, the
corresponding $\mu^{2}$ being its eigenvalues. A general fibre is both
stretched \emph{and} rotated, and the change in angle between two fibres is
the shear recorded by the off-diagonal $C_{AB}$. A rigid motion sends the
circle to a circle of equal radius, i.e.\ $\bC=\hat\I$.}
\end{figure}

\begin{remark}[reading Supplementary Figure B: why $\bF$ splits]
\label{rem:polarpreview}
The picture already contains the polar decomposition theorem of Chapter~3.
Turning the circle into the ellipse can be done in two stages: first stretch
along a fixed pair of orthogonal directions of $\kappa$, then rotate the
stretched figure into place. Symbolically $\bF=\bR\bU$ with $\bU$ symmetric
positive definite (the stretch) and $\bR$ orthogonal (the rotation).
Squaring kills the rotation, exactly as the picture suggests:
\[
   \bC=\bF^{t}\bF=(\bR\bU)^{t}(\bR\bU)=\bU^{t}\bR^{t}\bR\bU=\bU\hat\I\bU=\bU^{2},
\]
using $(\bB\bA)^{t}=\bA^{t}\bB^{t}$, $\bR^{t}\bR=\hat\I$ and
$\bU^{t}=\bU$. So $\bC$ records the ellipse's shape and size but nothing
about its final orientation, which is the precise sense in which it measures
distortion alone; $\mu^2=\ip\bM{\bC\bM}=\ip\bM{\bU^2\bM}=|\bU\bM|^2$ says the
stretch is done entirely by $\bU$. Symmetrically
$\bB=\bF\bF^{t}=\bR\bU^{2}\bR^{t}=\bV^{2}$ with $\bV=\bR\bU\bR^{t}$, the same
stretch seen from the current configuration.
\end{remark}

\begin{definition}[Lagrange and Euler strains]\label{def:strains}
The deviation of the stretch from its undeformed value $1$ is measured by
\begin{equation}
   \tfrac12(\mu^{2}-1)=\tfrac12\big(\ip\bM{\bC\bM}-1\big)=\ip\bM{\bE\bM},
   \qquad \bE=\tfrac12(\bC-\hat\I), \tag{2.96--2.97}\label{eq:2.96}
\end{equation}
the \emph{Lagrange strain tensor}, and by
\begin{equation}
   \tfrac12(1-\mu^{-2})=\ip\bmm{\vek{e}\bmm},
   \qquad \vek{e}=\tfrac12(\I-\bB^{-1}), \tag{2.98--2.99}\label{eq:2.98}
\end{equation}
the \emph{Euler strain tensor}. Both are symmetric.
\end{definition}

\begin{proof}[Verification]
Since $|\bM|=1$, $\ip\bM{\hat\I\bM}=\ip\bM\bM=1$, so
$\ip\bM{\bE\bM}=\tfrac12(\ip\bM{\bC\bM}-\ip\bM\bM)=\tfrac12(\mu^{2}-1)$,
which is \eqr{2.96}; the Euler case is identical using \eqr{2.91} and
$|\bmm|=1$.
\end{proof}

\begin{remark}[why the factor $\tfrac12$, and why strain at all]\label{rem:half}
For $\mu\simeq1$ write $\mu=1+\eps$ with $\eps$ small; then
$\mu^{2}-1=(1+\eps)^{2}-1=2\eps+\eps^{2}\simeq2(\mu-1)$. The factor $\tfrac12$
is chosen precisely so that $\ip\bM{\bE\bM}\simeq\mu-1$, the
\emph{extension} (change in length per unit original length), recovering the
engineer's strain in the small-deformation limit; likewise
$\ip\bmm{\vek{e}\bmm}\simeq\mu-1$. The reason for introducing strains at
all, rather than working with $\bC$ and $\bB$, is that they \emph{vanish} in
the undeformed state: by \eqr{2.79}, $\bF=\vek{1}$ gives
$\bC=\vek{1}^{t}\vek{1}=\hat\I$ and $\bB=\vek{1}\vek{1}^{t}=\I$
by \eqr{2.72}--\eqr{2.73}, hence $\bE=\hat\Ob$ and $\vek{e}=\Ob$. (Henceforth
we drop the caret on the referential zero tensor and let context decide.)
The names used here are not universal, and there are infinitely many
equivalent strain measures; see Ogden's book for a survey.
\end{remark}

\begin{remark}[symmetry is meaningless for two-point tensors]\label{rem:nosym}
$\bE$ and $\vek{e}$ are symmetric, but note carefully that symmetry and
skewness are defined only for tensors mapping a vector space \emph{to
itself}: $\bA^{t}=\pm\bA$ compares $\bA$ with an object of the same type.
For a two-point tensor such as $\bF$, the transpose $\bF^{t}$ maps $\Vt\to\hVt$
while $\bF$ maps $\hVt\to\Vt$, so the equation $\bF^{t}=\pm\bF$ cannot even be
written. Two-point tensors are therefore \emph{neither} symmetric nor skew,
which is another reason the deformation measures are built as $\bF^{t}\bF$
and $\bF\bF^{t}$.
\end{remark}

\paragraph*{Arclength of a deformed curve.}
On $C$, $\dl\bX(S)=\bM\,\dl S$ with $\dl S=|\dl\bX|$ (arclength
parametrisation), so the deformed arclength measure is
\begin{equation}
   \dl s=|\dl\bx|=\sqrt{\ip{\bF\dl\bX}{\bF\dl\bX}}
   =\sqrt{\ip\bM{\bC\bM}}\;\dl S=\mu(S,t_*)\,\dl S, \tag{2.100}\label{eq:2.100}
\end{equation}
that is $\dl s/\dl S=\mu(S,t_*)$ $(2.101)$, and the total length of the
deformed curve, and its change, are
\begin{equation}
   l(C_{t_*})=\int_{S_0}^{S_1}\!\mu(S,t_*)\,\dl S,\qquad
   l(C_{t_*})-l(C)=\int_{S_0}^{S_1}\!\big[\mu(S,t_*)-1\big]\,\dl S. \tag{2.102--2.103}\label{eq:2.102}
\end{equation}

\begin{figure}[htbp]\centering
\begin{tikzpicture}[scale=0.9]
  \draw[gray!75,thick] plot[smooth,tension=0.9] coordinates
    {(-1.35,-1.2) (-0.85,-0.25) (-0.15,0.4) (0.45,1.3) (1.1,2.05)};
  \fill (-1.35,-1.2) circle (1.7pt) node[left=2pt] {$S_0$};
  \fill (1.1,2.05) circle (1.7pt) node[right=2pt] {$S_1$};
  \fill (-0.15,0.4) circle (1.9pt); \node at (-0.6,0.55) {$S$};
  \draw[vec,black] (-0.15,0.4) -- ++(0.5,0.82);
  \node at (0.78,0.92) {$\bM$};
  \node at (0.35,-0.62) {\small$\dl\bX=\bM\,\dl S$};
  \node[gray!75] at (-1.55,1.45) {$C\subset\kappa$};
  \begin{scope}[shift={(5.9,0)}]
  \draw[gray!75,thick] plot[smooth,tension=0.9] coordinates
    {(-1.35,-1.2) (-0.55,-0.7) (-0.05,0.25) (0.35,1.25) (1.05,2.05)};
  \fill (-1.35,-1.2) circle (1.7pt) node[left=2pt] {$s=0$};
  \fill (1.05,2.05) circle (1.7pt) node[right=2pt] {$s=l(C_{t_*})$};
  \fill (-0.05,0.25) circle (1.9pt); \node at (-0.75,0.4) {$s(S,t_*)$};
  \node at (0.5,-0.62) {\small$\dl s=\mu\,\dl S$};
  \node[gray!75] at (-1.35,-1.75) {$C_{t_*}\subset\kappa_{t_*}$};
  \end{scope}
  \draw[mapar] (2.0,1.6) to[bend left=16] (3.75,1.6);
  \node at (2.9,2.2) {$\dfrac{\dl s}{\dl S}=\mu$};
\end{tikzpicture}
\caption{Arclength of a material point on material curves in $\kappa$ and $\kappa_{t_*}$.}
\label{fig:2.11}
\end{figure}

\begin{remark}[Figure 2.11]\label{rem:fig11}
Figure~2.11 shows $C\subset\kappa$ carrying the stations $S_0,S,S_1$, the
unit tangent $\bM$ and the element $\dl\bX=\bM\,\dl S$, beside its image
$C_{t_*}$ carrying the deformed arclength $s(S,t_*)$. The content of
$(2.101)$ is that the stretch $\mu$ is literally the local ratio of deformed
to undeformed arclength, so integrating it along the fibre gives the new
length: a one-dimensional version of $\dl v=J\,\dl V$.
\end{remark}

\paragraph*{Angles between material curves.}
Let two material curves cross at $\bX$ with unit tangents $\bM_1,\bM_2$ and
angle $\Theta$ ($\cos\Theta=\ip{\bM_1}{\bM_2}$); after deformation they cross
at $\bx$ with unit tangents $\bmm_1,\bmm_2$ and angle $\theta$
($\cos\theta=\ip{\bmm_1}{\bmm_2}$), where
\begin{equation}
   \mu_1\bmm_1=\bF\bM_1,\qquad \mu_2\bmm_2=\bF\bM_2 \tag{2.104}\label{eq:2.104}
\end{equation}
(Figure~2.12). Dotting the two and moving $\bF^{t}$ across,
\begin{equation}
   \mu_1\mu_2\cos\theta=\ip{\bF\bM_1}{\bF\bM_2}=\ip{\bM_1}{\bC\bM_2}, \tag{2.105}\label{eq:2.105}
\end{equation}
and dividing by $\mu_1\mu_2=\sqrt{\ip{\bM_1}{\bC\bM_1}}\sqrt{\ip{\bM_2}{\bC\bM_2}}$
from \eqr{2.83},
\begin{equation}
   \cos\theta=\frac{\ip{\bM_1}{\bC\bM_2}}
   {\sqrt{(\ip{\bM_1}{\bC\bM_1})(\ip{\bM_2}{\bC\bM_2})}} , \tag{2.106}\label{eq:2.106}
\end{equation}
in which $\ip{\bM_1}{\bC\bM_2}=\ip{\bM_2}{\bC\bM_1}$ because
$\ip{\bM_1}{\bC\bM_2}=\ip{\bC^{t}\bM_1}{\bM_2}=\ip{\bM_2}{\bC\bM_1}$ by the
symmetry of $\bC$. Running the same argument backwards with
$\mu_1^{-1}\bM_1=\bF^{-1}\bmm_1$, $\mu_2^{-1}\bM_2=\bF^{-1}\bmm_2$ from
\eqr{2.90} gives the undeformed angle in terms of current data,
\begin{equation}
   \cos\Theta=\frac{\ip{\bmm_1}{\bB^{-1}\bmm_2}}
   {\sqrt{(\ip{\bmm_1}{\bB^{-1}\bmm_1})(\ip{\bmm_2}{\bB^{-1}\bmm_2})}} . \tag{2.107--2.108}\label{eq:2.107}
\end{equation}

\begin{figure}[htbp]\centering
\begin{tikzpicture}[scale=0.85]
  % reference: two curves crossing at X
  \draw[thick] plot[smooth,tension=0.9] coordinates {(-1.7,-1.0) (-0.6,-0.35) (0,0) (0.9,0.9) (1.5,1.9)};
  \draw[thick] plot[smooth,tension=0.9] coordinates {(1.6,-0.9) (0.7,-0.5) (0,0) (-0.75,0.95) (-1.2,1.9)};
  \fill (0,0) circle (1.7pt) node[below=2pt] {$\bX$};
  \draw[vec] (0,0) -- (0.78,0.78); \node at (1.12,0.92) {$\bM_1$};
  \draw[vec] (0,0) -- (-0.62,0.85); \node at (-0.95,1.08) {$\bM_2$};
  \draw (0.42,0.42) arc (45:126:0.6); \node at (0.02,0.78) {$\Theta$};
  \node at (-1.55,1.55) {$\kappa$};
  % current
  \begin{scope}[shift={(6.0,0)}]
  \draw[thick] plot[smooth,tension=0.9] coordinates {(-1.5,-1.1) (-0.5,-0.45) (0,0) (0.75,1.05) (1.1,1.95)};
  \draw[thick] plot[smooth,tension=0.9] coordinates {(1.7,-0.7) (0.75,-0.4) (0,0) (-0.95,0.7) (-1.6,1.5)};
  \fill (0,0) circle (1.7pt) node[below=2pt] {$\bx$};
  \draw[vec] (0,0) -- (0.55,0.9); \node at (0.88,1.05) {$\bmm_1$};
  \draw[vec] (0,0) -- (-0.85,0.6); \node at (-1.22,0.75) {$\bmm_2$};
  \draw (0.3,0.49) arc (58:145:0.58); \node at (-0.12,0.82) {$\theta$};
  \node at (1.6,1.55) {$\kappa_t$};
  \end{scope}
  \draw[mapar] (2.15,1.35) to[bend left=16] (3.85,1.35);
  \node at (3.0,1.95) {$\bF$};
\end{tikzpicture}
\caption{Intersecting material curves in $\kappa$ and $\kappa_{t_*}$.}\label{fig:2.12}
\end{figure}

\begin{remark}[Figure 2.12, and why $\bC$ deserves the name ``distortion'']
\label{rem:fig12}
Figure~2.12 shows two material curves meeting at $\bX$ at angle $\Theta$ in
$\kappa$, and their images meeting at $\bx$ at angle $\theta$ in $\kappa_t$.
Formulas \eqr{2.106} and \eqr{2.108} say that $\bC$ (or $\bB$) determines
\emph{every} stretch and \emph{every} angle change at the point, and these
are the entire local geometric content of a deformation: hence the
Cauchy--Green tensors are said to characterise the \emph{state of distortion}
at $p$. Note the shear reading of the off-diagonal entries: taking
$\bM_1=\hbE_1$, $\bM_2=\hbE_2$ (initially orthogonal, $\Theta=\pi/2$),
\eqr{2.106} collapses to $\cos\theta=C_{12}/\sqrt{C_{11}C_{22}}$, so a nonzero
$C_{12}$ \emph{is} the loss of perpendicularity, while the diagonal entries
$C_{AA}=\mu^2$ (no sum) are squared stretches along the axes.
\end{remark}

%=======================================================================
\section{Problems}\label{sec:problems}
%=======================================================================

The three problems of \S2.8 are solved below in full index notation, using
only Chapter~1 algebra together with \eqr{2.39}, $F_{iA}=\chi_{i,A}$. Recall
the referential and spatial permutation symbols
$\eps_{ABC}=\tbox{\hbE_A}{\hbE_B}{\hbE_C}$ and
$\eps_{ijk}=\tbox{\be_i}{\be_j}{\be_k}$, and the contracted identities
$\eps_{ABC}\eps_{DBC}=2\dd_{AD}$, $\eps_{ijk}\eps_{ijk}=6$.

\begin{lemma}[symmetric contracted with antisymmetric vanishes]
\label{lem:symskw}
If $S_{BC}=S_{CB}$ and $A_{BC}=-A_{CB}$, then $A_{BC}S_{BC}=0$, the sum
running over both indices.
\end{lemma}

\begin{proof}
Use the antisymmetry of $A$, then relabel the two dummy indices, then use
the symmetry of $S$:
\[
   A_{BC}S_{BC}\ \overset{(1)}{=}\ -A_{CB}S_{BC}
   \ \overset{(2)}{=}\ -A_{BC}S_{CB}
   \ \overset{(3)}{=}\ -A_{BC}S_{BC}.
\]
Step (1) is $A_{BC}=-A_{CB}$; step (2) renames the summation indices; step
(3) is $S_{CB}=S_{BC}$. Hence $2A_{BC}S_{BC}=0$.
\end{proof}

The lemma is used twice below, at 2(a) and at 2(b).

\paragraph*{Problem 1: components of the cofactor of $\bF$.}
\emph{Show that} $F^{*}_{iA}=\tfrac12\eps_{ijk}\eps_{ABC}F_{jB}F_{kC}$.

\emph{Solution.} The cofactor is defined by $\bF^{*}(\ba\times\bb)
=\bF\ba\times\bF\bb$ for all $\ba,\bb\in\hVt$. Take $\ba=\hbE_B$,
$\bb=\hbE_C$, with
\[
   \hbE_B\times\hbE_C=\eps_{BCA}\hbE_A
\]
by Proposition~1.24. The left side becomes
\[
   \bF^{*}(\eps_{BCA}\hbE_A)=\eps_{BCA}\,\bF^{*}\hbE_A
   =\eps_{BCA}F^{*}_{iA}\,\be_i ,
\]
using $\bF^{*}\hbE_A=F^{*}_{iA}\be_i$ (the analogue of \eqr{2.40}). The right
side, expanding each image on $\{\be_i\}$ and using
$\be_j\times\be_k=\eps_{jki}\be_i$, is
\[
   \bF\hbE_B\times\bF\hbE_C=(F_{jB}\be_j)\times(F_{kC}\be_k)
   =F_{jB}F_{kC}\,\eps_{jki}\,\be_i .
\]
Equating components on $\be_i$ and writing $\eps_{jki}=\eps_{ijk}$ (cyclic),
\[
   \eps_{BCA}F^{*}_{iA}=\eps_{ijk}F_{jB}F_{kC}.
\]
Contract both sides with $\eps_{BCD}$ (sum over $B,C$). On the left,
$\eps_{BCA}\eps_{BCD}=2\dd_{AD}$, giving $2\dd_{AD}F^{*}_{iA}=2F^{*}_{iD}$;
on the right, $\eps_{BCD}=\eps_{DBC}$ (cyclic). Hence
\[
   F^{*}_{iD}=\tfrac12\,\eps_{ijk}\,\eps_{DBC}\,F_{jB}F_{kC},
\]
which upon renaming $D\to A$ is the required
$F^{*}_{iA}=\tfrac12\eps_{ijk}\eps_{ABC}F_{jB}F_{kC}$. \hfill$\square$

The identity $\eps_{BCA}\eps_{BCD}=2\dd_{AD}$ is Theorem~1.29. Directly:
for $A\neq D$ every term repeats an index and dies; for $A=D=3$ the
survivors are $\eps_{123}^{2}+\eps_{213}^{2}=2$. The factor $2$ is why
$\tfrac12$ appears, each unordered pair $\{B,C\}$ being counted twice.

\emph{Check.} For the simple shear of Chapter 3,
$F_{iA}=\dd_{iA}+\gamma\,\dd_{i1}\dd_{A2}$, whose matrix is
$\left(\begin{smallmatrix}1&\gamma&0\\0&1&0\\0&0&1\end{smallmatrix}\right)$.
Since $\det\bF=1$ and $\bF^{*}=J\bF^{-t}=\bF^{-t}$ by \eqr{2.55}, the answer
must be
$\left(\begin{smallmatrix}1&0&0\\-\gamma&1&0\\0&0&1\end{smallmatrix}\right)$.
Test the formula on the entry $A=1$, $i=2$:
\[
   F^{*}_{21}=\tfrac12\eps_{2jk}\eps_{1BC}F_{jB}F_{kC}.
\]
The only non-zero $\eps_{2jk}$ are $\eps_{231}=1$ and $\eps_{213}=-1$; the
only non-zero $\eps_{1BC}$ are $\eps_{123}=1$ and $\eps_{132}=-1$. So four
terms survive:
\[
\begin{aligned}
   2F^{*}_{21}&=\eps_{231}\eps_{123}F_{32}F_{13}
     +\eps_{231}\eps_{132}F_{33}F_{12}\\
   &\qquad+\eps_{213}\eps_{123}F_{12}F_{33}
     +\eps_{213}\eps_{132}F_{13}F_{32}\\
   &=(+1)(+1)(0)(0)+(+1)(-1)(1)(\gamma)\\
   &\qquad+(-1)(+1)(\gamma)(1)+(-1)(-1)(0)(0)\\
   &=-\gamma-\gamma=-2\gamma ,
\end{aligned}
\]
so $F^{*}_{21}=-\gamma$, as required.

\begin{figure}[htbp]\centering
\begin{tikzpicture}[scale=1.0,
  vec/.style={-{Stealth[length=2.2mm]},thick},
  lbl/.style={fill=white,inner sep=1pt,font=\small}]
% ---- reference parallelogram and its area vector ----
\begin{scope}
\fill[gray!14] (0,0)--(1.5,-0.35)--(2.1,0.95)--(0.6,1.30)--cycle;
\draw[thick] (0,0)--(1.5,-0.35)--(2.1,0.95)--(0.6,1.30)--cycle;
\draw[vec] (0,0)--(1.5,-0.35); \node[lbl,below] at (0.85,-0.30) {$\hbE_B$};
\draw[vec] (0,0)--(0.6,1.30);  \node[lbl,left]  at (0.52,1.05) {$\hbE_C$};
\draw[vec,gray!60] (1.05,0.475)--++(0.30,1.15);
\node[lbl] at (1.55,1.72) {$\hbE_B\times\hbE_C$};
\node[lbl] at (1.05,-1.50) {$\kappa$};
\end{scope}
% ---- arrow ----
\draw[-{Stealth[length=2.6mm]},thick,gray!70] (2.9,0.5)--(4.0,0.5);
\node[lbl,above,text=gray!85] at (3.45,0.56) {$\bF$};
% ---- image parallelogram ----
\begin{scope}[xshift=5.0cm]
\fill[gray!14] (0,0)--(2.1,0.15)--(2.5,1.15)--(0.4,1.0)--cycle;
\draw[thick] (0,0)--(2.1,0.15)--(2.5,1.15)--(0.4,1.0)--cycle;
\draw[vec] (0,0)--(2.1,0.15); \node[lbl,below] at (1.1,0.12) {$\bF\hbE_B$};
\draw[vec] (0,0)--(0.4,1.0);  \node[lbl,left]  at (0.34,0.82) {$\bF\hbE_C$};
\draw[vec,gray!60] (1.25,0.575)--++(0.18,1.25);
\node[lbl] at (1.75,1.92) {$\bF\hbE_B\times\bF\hbE_C$};
\node[lbl] at (1.25,-1.50) {$\kappa_{t}$};
\end{scope}
\draw[-{Stealth[length=2.6mm]},thick,gray!70,densely dashed]
   (1.35,1.72) .. controls (3.4,2.55) .. (5.9,1.95);
\node[lbl,font=\small] at (3.6,2.62) {$\bF^{*}$};
\end{tikzpicture}
\caption{Problem 1. $\bF$ carries the edge vectors of a material
parallelogram, $\bF^{*}$ carries its area vector. The problem converts the
implicit definition into components,
$F^{*}_{iA}=\tfrac12\eps_{ijk}\eps_{ABC}F_{jB}F_{kC}$; \eqr{2.55} then
converts those into $J\bF^{-t}$.}
\label{fig:cofactor}
\end{figure}

\paragraph*{Problem 2(a): $F^{*}_{iA}$ as a referential divergence.}
\emph{Show} $F^{*}_{iA}=G_{iAB,B}$, \emph{where}
$G_{iAB}=\tfrac12\eps_{ijk}\eps_{ABC}\chi_j\chi_{k,C}$.

\emph{Solution.} Differentiate $G_{iAB}$ with respect to $X_B$ and sum,
using the product rule (the $\eps$'s are constants):
\[
   G_{iAB,B}=\tfrac12\eps_{ijk}\eps_{ABC}
   \big(\chi_{j,B}\,\chi_{k,C}+\chi_j\,\chi_{k,CB}\big).
\]
Second term: for fixed $A,k$ put $S_{BC}=\chi_{k,CB}$ and
$A_{BC}=\eps_{ABC}$. Mixed partials commute, so $S_{BC}=S_{CB}$; $\eps$
changes sign under a swap, so $A_{BC}=-A_{CB}$. By
Lemma~\ref{lem:symskw},
\[
   \eps_{ABC}\,\chi_{k,CB}=0 .
\]
First term, by $F_{jB}=\chi_{j,B}$ and Problem~1,
\[
   \tfrac12\eps_{ijk}\eps_{ABC}F_{jB}F_{kC}=F^{*}_{iA}.
\]
Hence $G_{iAB,B}=F^{*}_{iA}$. \hfill$\square$

\paragraph*{Problem 2(b): the Piola identity.}
\emph{Show} $G_{iAB}=-G_{iBA}$ \emph{and deduce} $F^{*}_{iA,A}=0$.

\emph{Solution.} Swapping the free indices $A,B$ in the definition and using
only the antisymmetry $\eps_{BAC}=-\eps_{ABC}$,
\[
   G_{iBA}=\tfrac12\eps_{ijk}\eps_{BAC}\chi_j\chi_{k,C}
   =-\tfrac12\eps_{ijk}\eps_{ABC}\chi_j\chi_{k,C}=-G_{iAB}.
\]
Now take the referential divergence of Problem 2(a):
\[
   F^{*}_{iA,A}=(G_{iAB,B})_{,A}=G_{iAB,BA}.
\]
Apply Lemma~\ref{lem:symskw} again on the pair $(A,B)$: $(\cdot)_{,BA}$ is
symmetric, $G_{iAB}$ antisymmetric. Hence
\[
   F^{*}_{iA,A}=0,\qquad\text{i.e.}\qquad \Dvg\bF^{*}=\bze,
\]
the \emph{Piola identity}. \hfill$\square$

\begin{figure}[htbp]\centering
\begin{tikzpicture}[scale=1.0,
  vec/.style={-{Stealth[length=2.1mm]},semithick},
  lbl/.style={fill=white,inner sep=1pt,font=\small}]
% ---- closed material surface in kappa ----
\begin{scope}
\draw[thick] plot[smooth cycle,tension=0.85] coordinates
  {(0,1.05) (1.05,0.66) (1.35,-0.32) (0.55,-1.05) (-0.68,-0.86) (-1.18,0.05)};
\foreach \a/\b/\c/\d in {0/1.05/0.10/0.55, 1.05/0.66/0.60/0.28,
  1.35/-0.32/0.62/-0.18, 0.55/-1.05/0.18/-0.56, -0.68/-0.86/-0.42/-0.48,
  -1.18/0.05/-0.60/0.06}
  \draw[vec,gray!75] (\a,\b)--++(\c,\d);
\node[lbl] at (0,-1.60) {$\partial P\subset\kappa$};
\node[lbl,font=\small] at (1.45,0.95) {$\bN$};
\end{scope}
% ---- arrow ----
\draw[-{Stealth[length=2.6mm]},thick,gray!70] (2.0,0)--(3.1,0);
\node[lbl,above,text=gray!85] at (2.55,0.06) {$\bchi_{\kappa}$};
% ---- image in kappa_t ----
\begin{scope}[xshift=4.9cm,rotate=-18,xscale=1.25,yscale=0.85]
\draw[thick] plot[smooth cycle,tension=0.85] coordinates
  {(0,1.05) (1.05,0.66) (1.35,-0.32) (0.55,-1.05) (-0.68,-0.86) (-1.18,0.05)};
\foreach \a/\b/\c/\d in {0/1.05/0.10/0.55, 1.05/0.66/0.60/0.28,
  1.35/-0.32/0.62/-0.18, 0.55/-1.05/0.18/-0.56, -0.68/-0.86/-0.42/-0.48,
  -1.18/0.05/-0.60/0.06}
  \draw[vec,gray!75] (\a,\b)--++(\c,\d);
\end{scope}
\node[lbl] at (4.9,-1.60) {$\partial P_{t}\subset\kappa_{t}$};
\node[lbl,font=\small] at (6.5,1.05) {$\bn$};
% ---- the statement ----
\node[align=center,font=\small] at (2.9,-2.55)
  {$\displaystyle\int_{\partial P}\bF^{*}\bN\,\dl A
    =\int_{P}\Dvg\bF^{*}\,\dl V=\bze
    \qquad\Longleftrightarrow\qquad
    \int_{\partial P_{t}}\bn\,\dl a=\bze$};
\end{tikzpicture}
\caption{The Piola identity. A closed material surface stays closed, so its
normals integrate to zero in $\kappa_{t}$ as they do in $\kappa$. By
Piola--Nanson \eqr{2.54} that is the divergence theorem applied to
$\bF^{*}$, and it holds for every $P$ and every motion only if
$\Dvg\bF^{*}=\bze$.}
\label{fig:piola}
\end{figure}

\begin{remark}[what the Piola identity buys]\label{rem:piolauses}
$\Dvg\bF^{*}=\bze$ says the cofactor field is automatically divergence free,
no matter what motion is prescribed. Combined with Piola--Nanson
\eqr{2.54}, it is exactly the statement that the deformed images of closed
material surfaces stay closed: applying the divergence theorem,
$\int_{\partial P}\bF^{*}\bN\,\dl A=\int_P\Dvg\bF^{*}\,\dl V=\bze$, i.e.\
$\int_{\partial P_t}\bn\,\dl a=\bze$, the elementary fact that the outward
normals of any closed surface integrate to zero. It is the kinematic
identity that later converts the Cauchy stress balance into its Piola form.
\end{remark}

\paragraph*{Problem 2(c): $\det\bF$ as a divergence.}
\emph{Show} $\det\bF=\tfrac13F_{iA}F^{*}_{iA}$, \emph{and hence that there is
a field $\vek{H}(\bX,t)=H_A\hbE_A$ with} $\det\bF=H_{A,A}$.

\emph{Solution.} The cofactor identity $\bA^{t}\bA^{*}=(\det\bA)\I$ of
Problem 22(c) of Chapter~1 holds verbatim for the two-point $\bF$, where it
reads, in components,
\[
   F_{iA}F^{*}_{iB}=(\det\bF)\,\dd_{AB}.
\]
Setting $B=A$ and summing, with $\dd_{AA}=3$,
\[
   F_{iA}F^{*}_{iA}=3\det\bF
   \qquad\Longrightarrow\qquad \det\bF=\tfrac13F_{iA}F^{*}_{iA}.
\]
Now substitute $F_{iA}=\chi_{i,A}$ and use the Piola identity of part (b) to
move the derivative outside:
\[
   \det\bF=\tfrac13\chi_{i,A}F^{*}_{iA}
   =\tfrac13\Big[(\chi_iF^{*}_{iA})_{,A}-\chi_i\underbrace{F^{*}_{iA,A}}_{=0}\Big]
   =\Big(\tfrac13\chi_iF^{*}_{iA}\Big)_{,A}.
\]
Hence $\det\bF=H_{A,A}$ with
\[
   H_A=\tfrac13\,\chi_i\,F^{*}_{iA},
\]
i.e.\ $\vek{H}=\tfrac13\bF^{*t}\bchi_\kappa$. \hfill$\square$

\begin{remark}[why one wants $J$ as a divergence]\label{rem:Hdiv}
Writing $J=\det\bF=\Dvg\vek{H}$ turns the volume integral of $J$ into a
surface integral by the divergence theorem,
$\int_P J\,\dl V=\int_{\partial P}\ip{\vek{H}}{\bN}\,\dl A$. Since
$\int_PJ\,\dl V=\int_{P_t}\dl v$ is the deformed volume, this expresses the
volume enclosed by a deformed surface purely in terms of data \emph{on} that
surface, the familiar ``volume from a surface integral'' trick, and it is the
kinematic origin of several variational identities in nonlinear elasticity.
\end{remark}

\paragraph*{Problem 2(d): $\det\bF^{t}=\det\bF$.}

\emph{Solution.} Substituting Problem 1 into the result of 2(c),
\begin{equation}
   \det\bF=\tfrac13F_{iA}\cdot\tfrac12\eps_{ijk}\eps_{ABC}F_{jB}F_{kC}
   =\tfrac16\,\eps_{ijk}\,\eps_{ABC}\,F_{iA}F_{jB}F_{kC},
   \tag{$\ast$}
\end{equation}
the two-point version of the contracted determinant formula of Chapter~1.
The transpose has components $(\bF^{t})_{Ai}=F_{iA}$, and $\bF^{t}$ maps
$\Vt\to\hVt$, so the same formula with the roles of the two index families
exchanged gives
\[
\begin{aligned}
   \det\bF^{t}&=\tfrac16\,\eps_{ABC}\,\eps_{ijk}\,
      (\bF^{t})_{Ai}(\bF^{t})_{Bj}(\bF^{t})_{Ck}\\
   &=\tfrac16\,\eps_{ABC}\,\eps_{ijk}\,F_{iA}F_{jB}F_{kC}=\det\bF,
\end{aligned}
\]
since the two scalar expressions differ only in the order in which the same
factors are written. \hfill$\square$

\paragraph*{Problem 3: change of basis for a two-point tensor.}
\emph{With $\{\hbE_A\},\{\hbE'_A\}$ orthonormal in $\hVt$ and
$\{\be_i\},\{\be'_i\}$ orthonormal in $\Vt$, and
$\bF=F_{iA}\be_i\otimes\hbE_A=F'_{iA}\be'_i\otimes\hbE'_A$, find $F'_{iA}$.}

\emph{Solution.} Introduce the two independent sets of direction cosines,
one for each space,
\[
   a_{ij}=\ip{\be_i}{\be'_j}\ \ (\text{spatial}),\qquad
   \hat a_{AB}=\ip{\hbE_A}{\hbE'_B}\ \ (\text{referential}),
\]
so that, by the orthonormal expansion of Chapter~1,
\[
   \be'_i=(\ip{\be_j}{\be'_i})\be_j=a_{ji}\be_j,\qquad
   \hbE'_A=(\ip{\hbE_B}{\hbE'_A})\hbE_B=\hat a_{BA}\hbE_B .
\]
Extract the primed components with $(2.41)$, $F'_{iA}=\ip{\be'_i}{\bF\hbE'_A}$.
First act with $\bF$ on the primed referential basis vector, using linearity
and $\bF\hbE_B=F_{jB}\be_j$ from \eqr{2.40}:
\[
   \bF\hbE'_A=\hat a_{BA}\,\bF\hbE_B=\hat a_{BA}F_{jB}\,\be_j .
\]
Then dot with $\be'_i=a_{ki}\be_k$ and use $\ip{\be_k}{\be_j}=\dd_{kj}$:
\[
   F'_{iA}=\ip{a_{ki}\be_k}{\hat a_{BA}F_{jB}\be_j}
   =a_{ki}\hat a_{BA}F_{jB}\,\dd_{kj}
   =a_{ji}\,\hat a_{BA}\,F_{jB}.
\]
Hence
\[
   \boxed{\;F'_{iA}=a_{ji}\,\hat a_{BA}\,F_{jB}\;}
\]
\hfill$\square$

\emph{In matrix form.} With $(a)_{ij}=a_{ij}$ and
$(\hat a)_{AB}=\hat a_{AB}$, the index pattern $a_{ji}(\cdot)_{jB}
\hat a_{BA}$ is
\[
   (F')=(a)^{t}\,(F)\,(\hat a).
\]
Both $(a)$ and $(\hat a)$ are orthogonal, by the computation of Problem
9(a) of Chapter 1 applied in each space separately:
\[
   a_{ki}a_{kj}=\dd_{ij},\qquad \hat a_{CA}\hat a_{CB}=\dd_{AB}.
\]

\emph{Check 1: $J$ is unchanged.} Taking determinants,
\[
   \det(F')=\det(a)^{t}\det(F)\det(\hat a)=(\pm1)\det(F)(\pm1),
\]
and both signs are $+1$ when the two changes of basis preserve
handedness. So $J$ is basis independent, as \eqr{2.50} requires.

\emph{Check 2: $\bC$ transforms with one matrix only.} From
$C_{AB}=F_{iA}F_{iB}$,
\[
\begin{aligned}
   C'_{AB}&=F'_{iA}F'_{iB}
   =\big(a_{ji}\hat a_{CA}F_{jC}\big)\big(a_{ki}\hat a_{DB}F_{kD}\big)\\
   &=\hat a_{CA}\hat a_{DB}\,\dd_{jk}F_{jC}F_{kD}
   =\hat a_{CA}\hat a_{DB}\,C_{CD},
\end{aligned}
\]
using $a_{ji}a_{ki}=\dd_{jk}$. The spatial cosines have cancelled: $\bC$
is a referential tensor and cannot feel a change of spatial basis. The
same computation on $B_{ij}=F_{iA}F_{jA}$ cancels the referential cosines
instead.

\begin{remark}[the transformation law identifies the tensor's type]
\label{rem:2ptlaw}
Compare with Chapter~1: a vector picks up one factor of the cosine matrix,
$v'_i=a_{ki}v_k$; an ordinary (one-point) tensor picks up two factors of the
\emph{same} matrix, $A'_{ij}=a_{ki}a_{lj}A_{kl}$. A two-point tensor also
picks up two factors, but one from \emph{each} space, $a_{ji}$ for the
lower-case index and $\hat a_{BA}$ for the upper-case one. The two rotations
are independent: one may re-label the reference frame without touching the
current frame, and $\bF$ responds only in its capital index. This is the
precise sense in which $\bF$ ``has one foot in each configuration'', and it
is why the alphabet convention of Remark~\ref{rem:fig6} is more than
cosmetic.
\end{remark}

%=======================================================================
\appendix
\chapter{Dictionary to the book}\label{app:dict}
%=======================================================================

\renewcommand{\arraystretch}{1.2}
\begin{center}\scriptsize
\begin{tabular}{@{}ll@{\hspace{1.6em}}ll@{}}
\hline
Book & Here & Book & Here\\
\hline
\eqr{2.1} motion & Def.~\ref{def:motion}
 & \eqr{2.34}--\eqr{2.37} & \S\ref{sec:defgrad}, Thm.~\ref{thm:GradgradF}\\
\eqr{2.2} velocity, accel. & Def.~\ref{def:va}
 & \eqr{2.38}--(2.41) & eqs.\ here in \S\ref{sec:defgrad}\\
\eqr{2.3}, \eqr{2.4} reference & Def.~\ref{def:ref}
 & \eqr{2.42}--\eqr{2.44} & Lem.~\ref{lem:dyadprod}\\
\eqr{2.5}--\eqr{2.7} $\bchi_\kappa$ & eqs.\ in \S\ref{sec:pva}
 & \eqr{2.45}, \eqr{2.46} & material curves, \S\ref{sec:defgrad}\\
\eqr{2.8}--\eqr{2.10} velocity forms & \S\ref{sec:pva}
 & \eqr{2.47}--\eqr{2.50} volume & eqs.\ here, Rem.~\ref{rem:fig8}\\
\eqr{2.11}--\eqr{2.13} descriptions & \S\ref{sec:desc}, Rem.~\ref{rem:decorations}
 & \eqr{2.51}--\eqr{2.54} Piola--Nanson & \S\ref{sec:defgrad}\\
\eqr{2.14}--\eqr{2.16} differentials & \S\ref{sec:desc}
 & \eqr{2.55} $\bF^{*}=J\bF^{-t}$ & \S\ref{sec:defgrad}\\
\eqr{2.17}--\eqr{2.21} material deriv. & Thm.~\ref{thm:matder}
 & \eqr{2.56}--\eqr{2.59} mass & \S\ref{sec:mass}\\
\eqr{2.22}--\eqr{2.24} coordinates & \S\ref{sec:cart}, Rem.~\ref{rem:fig6}
 & \eqr{2.60}--\eqr{2.65} localization & Thm.~\ref{thm:localization}\\
\eqr{2.25}--\eqr{2.29} components & \S\ref{sec:cart}
 & Example \eqr{2.30}--\eqr{2.33} & Example~\ref{ex:motion}\\
\eqr{2.66}, \eqr{2.67} stretch & \S\ref{sec:shifter}, Rem.~\ref{rem:fig10}
 & \eqr{2.89}--\eqr{2.95} $\bB$ & Prop.~\ref{prop:B}\\
\eqr{2.68}--\eqr{2.74} shifter & Def.~\ref{def:shifter}, Prop.~\ref{prop:shifter}
 & \eqr{2.96}--\eqr{2.99} strains & Def.~\ref{def:strains}, Rem.~\ref{rem:half}\\
\eqr{2.75}--\eqr{2.78} displacement & \S\ref{sec:shifter}, Rem.~\ref{rem:shifteruse}
 & \eqr{2.100}--\eqr{2.103} arclength & \S\ref{sec:cauchygreen}, Rem.~\ref{rem:fig11}\\
\eqr{2.79}--\eqr{2.81} undeformed & Prop.~\ref{prop:undeformed}
 & \eqr{2.104}--\eqr{2.108} angles & \S\ref{sec:cauchygreen}, Rem.~\ref{rem:fig12}\\
\eqr{2.82}--\eqr{2.88} $\bC$ & Prop.~\ref{prop:C}
 & \S2.8 Problems 1--3 & \S\ref{sec:problems}\\
\hline
\end{tabular}
\end{center}

\medskip\noindent
Figures 2.1--2.12 are described in Remarks~\ref{rem:fig123}, \ref{rem:fig4},
\ref{rem:fig5}, \ref{rem:fig6}, \ref{rem:fig7}, \ref{rem:fig8},
\ref{rem:fig9}, \ref{rem:fig10}, \ref{rem:fig11} and \ref{rem:fig12}. All
algebra uses only Chapter~1: the tensor product and its component rule, the
dyad product $(\ba\otimes\bb)(\bc\otimes\bd)=(\ip\bb\bc)\ba\otimes\bd$, the
transpose rules of Problems 17--20, the box product and determinant
$\tbox{\bF\ba}{\bF\bb}{\bF\bc}=\det\bF\,\tbox\ba\bb\bc$, the cofactor
identity $\bA^{t}\bA^{*}=(\det\bA)\I$ of Problem 22(c), the contracted
$\eps$--$\dd$ identities, and the coordinate free chain rule
$\dl f=\ip{\grad f}{\dl\bx}$.

\end{document}
